\section{Involutions of the three-cusp group}

Throughout this part \(F=\widehat F(x,y)\), \(z=(xy)^{-1}\), and
\(\mathcal P\) denotes the subgroup of \(\operatorname{Out}(F)\)
preserving each named peripheral conjugacy class up to a common
unit power. All group homomorphisms are continuous. We use the
normalized lift
\[
 (\lambda,f):\quad x\longmapsto x^\lambda,\qquad
 y\longmapsto f^{-1}y^\lambda f,\qquad f\in F',
\]
where \(F'\) is the closed derived subgroup. This normalization is
a homomorphic section, as in the refereed LM4 D.32. Write
\(c=(-1,1)\). We use R1--R3 only with their stated scopes.

\statement{A.1}{Known}{Torsion in an acting group and its action on torsion are different assertions}
Lochak and Schneps (1997, Proposition 4(i)) prove
\(C_F(c)=1\); Proposition 4(ii) proves that \(c\) is
self-centralizing in \(\widehat{GT}\).
Schneps (1997, \S1.4(4)) and Lochak and Schneps (2006,
Question 1.4) ask whether all nonidentity torsion in
\(\widehat{GT}\) consists of conjugates of \(c\).
Harbater and Schneps (2000, Main Theorem(a), proved by
Propositions 3--4 and Theorem 2) identify
\(\widehat{GT}_0\) with the centralizer in \(\mathcal P\) of
the geometric \(S_3\).
Collas's genus-zero Theorems 3.3 and 3.10 identify prime-order
torsion in \(\widehat\Gamma_{0,[n]}\) with discrete torsion
and prove the \(\lambda\)-conjugacy action of \(\widehat{GT}\).
His genus-one Theorems 6.3 and 6.5 give the corresponding
statements with acting group \(\widehat{GR}\), isomorphic
to the subgroup denoted \(I\Gamma\) in that article;
replacing it there by all of \(\widehat{GT}\) is not
justified by those theorems.\cite{A_LS1997,A_Schneps1997,A_LS2006,A_HS2000,A_Collas2012a,A_Collas2012b}

\begin{proof}
For the centralizer comparison identify \(K(0,4)\) with
\(\widehat F(x,y)\). The natural \(S_4\)-action on its outer
automorphism group factors through the anharmonic \(S_3\).
Harbater--Schneps's quasi-special condition is precisely the
three peripheral power condition here; their symmetric
condition is commutation with that \(S_3\). Their proof
translates commutation with the transposition into the
two-cycle relation and commutation with the three-cycle into
the hexagon. The identification has the required source
and target, not merely an abstract isomorphism.

In Collas's formulas the variable of prime order lies in a
completed mapping class group; \((\lambda,f)\) is an arbitrary
element of the acting group. The conclusion
\(F(\gamma)=u\gamma^\lambda u^{-1}\) does not classify
finite-order elements \(F\) of the acting group. The subgroup
\(\widehat{GR}\) in the genus-one article is retained in its
Theorem 6.5. Neither action theorem therefore proves task A.8.
\end{proof}

The opened sources establish no literature antecedent for the
full order bound in R1. Searches included the authors' source
lists, the cited torsion papers, and recent work on profinite
curve-complex automorphisms and Galois--Teichm\"uller theory.
This records the search, not a priority claim or a proof that
no antecedent exists.

\statement{A.2}{Proved}{Relative pronilpotent levels recover the full peripheral group}
For an open characteristic subgroup \(U\) of \(F\), put
\[
 R_U=R_{\mathrm{nil}}(U),\quad F_U=F/R_U,\quad
 P_U=\operatorname{Out}^{\sharp}(F_U).
\]
For \(V\subset U\), the quotient \(F_V\to F_U\) induces
a continuous epimorphism \(P_V\to P_U\). There is an
isomorphism of profinite groups
\[
                 \mathcal P \xrightarrow{\sim}\varprojlim_U P_U.
\]
Two elements of \(\mathcal P\) are conjugate in \(\mathcal P\)
if and only if their images are conjugate in every \(P_U\).
This proves Lead 1.

\begin{proof}
Restricting finite nilpotent quotients of \(U\) to \(V\)
shows \(R_V\subset R_U\). Each \(R_U\) is characteristic in
\(F\). Since \(R_U\subset U\) and open characteristic
subgroups are cofinal in the finitely generated profinite
group \(F\), we have \(\bigcap_U R_U=1\).

We first prove surjectivity onto \(P_U\). Choose a representative
of a peripheral outer class with common multiplier \(\lambda\),
and a Galois element with cyclotomic multiplier \(\lambda\);
the cyclotomic character is onto
\(\widehat{\mathbb Z}^{\times}\). The latter acts on \(F_U\).
Composing with its inverse gives a multiplier-one peripheral
automorphism. R2 lifts this automorphism, with chosen
conjugators, to \(F\). Composing back proves surjectivity.
In particular every peripheral automorphism of \(F_V\)
lifts to \(F\), so preserves the image of \(R_U\).
This proves that the transition maps are defined on the
entire displayed outer groups. Their surjectivity follows
by lifting first to \(\mathcal P\).

Here the automorphism groups are compact: intersecting
kernels of all maps to finite groups of bounded order
gives a cofinal system of open fully invariant subgroups.
The conditions that a continuous endomorphism induces
bijections on these finite quotients are closed conditions
on its generator images. Consequently \(\operatorname{Aut}(F)\)
is profinite. Inner automorphisms form a compact subgroup.
The peripheral conditions are closed, since conjugators
and the common multiplier range over compact groups.
The same argument applies to the finitely generated
quotients \(F_U\).

If \([a]\in\mathcal P\) dies in every \(P_U\), the sets
\[
 \{g\in F:a(x)\equiv gxg^{-1},\
              a(y)\equiv gyg^{-1}\pmod {R_U}\}
\]
are nonempty compact sets with the finite intersection
property. Their intersection gives an actual inner
representative for \(a\), proving injectivity.
For a compatible family of classes in \(P_U\), their
nonempty compact inverse images in \(\mathcal P\) have
the same finite intersection property, proving surjectivity.
Finally apply this argument to the transporter sets
\(\{g:g_Ua_Ug_U^{-1}=b_U\}\) in \(\mathcal P\).
Surjectivity onto \(P_U\) makes these nonempty whenever
the two images are conjugate. Their intersection is the
required full conjugator.
\end{proof}

\statement{A.3}{Proved}{A central quotient is the exact obstruction to a peripheral lift}
Let \(N\) be a closed characteristic subgroup of \(F\), and
let \([N,F]\) denote the closed subgroup generated by
commutators between \(N\) and \(F\).
An automorphism \(a\) of \(F/N\), with specified conjugators
\(b_i\) satisfying
\[
 a(\bar x_i)=b_i\bar x_i b_i^{-1},\qquad
 (x_1,x_2,x_3)=(x,y,z),
\]
lifts to a peripheral automorphism of \(F\) with these
conjugators if and only if it lifts to a peripheral
\emph{automorphism} of \(F/[N,F]\) with conjugators above
the \(b_i\). This proves Lead 2.

\begin{proof}
Necessity follows from characteristicity. For sufficiency
let \(a^\#\) and \(b_i^\#\) be the data on \(F/[N,F]\).
We use the finite central-refinement count in the refereed
LM7\_PUSH D.50, equations (D.50.6)--(D.50.16).
For any actual quotient \(q_E:F\twoheadrightarrow E\)
and normal subgroup \(A\triangleleft E\), that count
constructs an actual quotient
\(q^\dagger:F\twoheadrightarrow D^\dagger\), with central
kernel over \(D=E/A\). A second epimorphism to
\(D^\dagger\) with specified peripheral conjugators has
the same number of epimorphism lifts to \(E\), with the
transported \(A\)-conjugacy classes on the boundary.
This counting assertion uses only normality of \(A\).
The hypothesis \([E,A]=A\) occurs in the separate
joint-image assertion (D.50.17), not used here.

Take \(A=q_E(N)\). The image \(q^\dagger(N)\) is central,
so \(q^\dagger\) kills \([N,F]\).
Thus \(q^\dagger a^\#\) is an epimorphism from the
same source to \(D^\dagger\), with conjugators the
images of \(b_i^\#\). The finite count gives an epimorphism
to \(E\) whose conjugators agree with the \(b_i\) in \(D\);
the reference count is positive since it includes \(q_E\).

In the compact space \(F^3\) of conjugators impose the
boundary product relation, generation in every finite
quotient, and agreement with each \(b_i\) in every finite
quotient of \(F/N\). A finite list of conditions lies in
one joint finite quotient \(E\), and the preceding
construction satisfies that list. Compactness gives
full conjugators satisfying all conditions. Their
endomorphism is surjective since its compact image
maps onto every finite quotient; Hopficity makes it
an automorphism. Its generator images induce exactly
\(a\) modulo \(N\).
\end{proof}

\statement{A.4}{Proved}{A degree-three coefficient imposed by the third cusp}
Over a field \(k\) of characteristic zero or greater
than three, use the free class-three Lie algebra on
\(X,Y\), and set
\[
 Z=[X,Y],\quad U_3=[X,Z],\quad V_3=[Y,Z],\quad
 x=\exp X,\quad y=\exp Y .
\]
If \(x\mapsto x^\mu,\ y\mapsto g^{-1}y^\mu g\), with
\(\mu\in k^\times\) and
\(\log g=tZ+\text{central terms}\), sends \(xy\)
to a conjugate of \((xy)^\mu\), then
\(t=(\mu^2-1)/24\).

\begin{proof}
The Baker--Campbell--Hausdorff formula gives
\[
 W=\log(xy)=X+Y+\tfrac12Z+\tfrac1{12}(U_3-V_3).
\]
The logarithm of the image is
\[
 \mu(X+Y)+\tfrac{\mu^2}{2}Z+
       \tfrac{\mu^3}{12}(U_3-V_3)+\mu tV_3 .
\]
Write the logarithm of a conjugator as
\(aX+bY+dZ+\text{central terms}\).
In \(\exp(\operatorname{ad}(aX+bY+dZ))(\mu W)\),
comparison of degree two gives \(m=a-b=(\mu-1)/2\).
The difference of the \(U_3\)- and \(V_3\)-coefficients,
divided by \(\mu\), then gives
\[
 \frac{\mu^2}{6}-t=\frac16+\frac m2+\frac{m^2}{2}
                  =\frac{3\mu^2+1}{24}.
\]
The result follows. The coefficient \(d\) and all
central coordinates cancel.
\end{proof}

\statement{A.5}{Refuted}{The specified conjugator in task A.3 need not preserve the third cusp}
There is an involution \(\alpha\in\widehat{GT}\), conjugate
to \(c\) by a Galois element, for which the unique solution
\(h\in F'\) of \(f_\alpha=h^{-1}\alpha(h)\) gives
\[
                 \psi(x)=x,\qquad\psi(y)=hyh^{-1}
\]
that does not preserve the conjugacy class of \(z\)
on \(F/R_{\mathrm{nil}}(F)\). The failure occurs in
the class-three exponent-seven quotient of order \(7^5\).

\begin{proof}
Choose \(\theta=(\mu,g)\) in the Galois image with
\(\mu\equiv2\pmod7\), possible by surjectivity of the
cyclotomic character modulo seven, and put
\(\alpha=\theta c\theta^{-1}\).
Normalization is a homomorphic section, so these
are actual automorphisms with that equality.
From \(\alpha\theta(y)=\theta(y)^{-1}\) we obtain
\[
 \alpha(y)=\alpha(g)g^{-1}y^{-1}g\alpha(g)^{-1},
 \qquad f_\alpha=g\alpha(g)^{-1}.
\]
Thus \(h=g^{-1}\) is a solution. Also
\(C_F(\alpha)=\theta(C_F(c))=1\) by
Lochak--Schneps Proposition 4(i). The quotient of two
solutions is fixed by \(\alpha\), so \(h\) is unique.

Let \(Q=\exp L\) for the free class-three Lie algebra
\(L\) on two generators over \(\mathbb F_7\).
It has order \(7^5\), exponent seven, and is generated
by \(x=\exp X,y=\exp Y\); hence it is a quotient of
\(F^{(7)}\) and of \(F/R_{\mathrm{nil}}(F)\).
The third-cusp condition on \(\theta\) and A.4 give
\[
       \operatorname{coeff}_Z(\log\bar g)
                    =\frac{2^2-1}{24}=1 .
\]
But \(\psi\) has multiplier one and the same conjugator
\(g\). Peripheral preservation for \(\psi\) would force
this coefficient to be zero by A.4, a contradiction.

The exact certificate \texttt{check\_A3\_class3.py}
realizes \(Q\) in the units of
\[
 \mathbb F_7\langle X,Y\rangle/(\text{words of length }4).
\]
It enumerates all \(16807\) elements; the conjugacy
orbit of \(xy\) has \(49\) elements. For
\(g=\exp[X,Y]\) it verifies
\[
 xg^{-1}yg\not\sim xy,\qquad
 x^2g^{-1}y^2g\sim(xy)^2 .
\]
The actual Galois conjugator differs only by central
degree-three coordinates relevant to this calculation,
which do not affect it. This certifies a failure in
an actual full profinite example; no order-preserving
lift of a finite involution is being assumed.
\end{proof}

The correct universal statement from R1 is that this
\(\psi\) intertwines \(\alpha,c\) on the maximal
pro-\(p\) quotients as automorphisms of free groups.
Peripheral preservation of this specific \(\psi\)
is an additional condition and is false.

\statement{A.6}{Proved}{The pure centralizer of standard conjugation has two elements}
\[
 C_{\mathcal P}(c)=N_{\mathcal P}(\langle c\rangle)=\{1,c\}.
\]
If \(\alpha\in\widehat{GT}_0\), every conjugator in
\(\mathcal P\) from \(c\) to \(\alpha\) belongs
to \(\widehat{GT}_0\).

\begin{proof}
A normalized \((\mu,g)\) commuting with \(c\) satisfies
\(c(g)=g\), by the normalized group law. A.1 gives
\(g=1\). In particular the multiplier-one centralizer
is trivial. Let \(\Sigma=S_3\) be the geometric
permutation group. If
\(\beta c\beta^{-1}=\alpha\) and \(\alpha,c\) commute
with \(\Sigma\), then
\(r_\sigma=\beta^{-1}\sigma\beta\sigma^{-1}\)
is pure, has multiplier one, and centralizes \(c\).
Thus \(r_\sigma=1\), and the Harbater--Schneps
identification in A.1 puts \(\beta\) in
\(\widehat{GT}_0\).

Apply this to a centralizing \((\mu,1)\).
The hexagon, with \(m=(\mu-1)/2\), reads
\((xy)^m=y^mx^m\). In the finite Heisenberg groups
with multiplication
\((a,b,d)(a',b',d')=(a+a',b+b',d+d'+ab')\),
the final coordinates of these two words are
\(m(m+1)/2\) and zero. Hence \(\mu^2=1\) in
\(\widehat{\mathbb Z}\). The actual automorphism
\((\mu,1)\) then has square one. R1 forces its
multiplier to be globally \(1\) or \(-1\), proving
the centralizer assertion. A normalizer of
\(\langle c\rangle\) centralizes its unique
nonidentity element.
\end{proof}

\statement{A.7}{Proved}{The multiplier-one cohomology set already has continuum many classes}
For \(G=\mathcal P,\widehat{GT}_0\), or \(\widehat{GT}\),
put \(H_G=\ker\lambda\), with \(C_2\) acting by \(c\).
There is a canonical closed injection of profinite spaces
\[
 \widehat{\mathbb Z}^{\times}/\{1,-1\}
 \longrightarrow H^1(C_2,H_G),\qquad
 \mu\longmapsto[\theta_\mu c\theta_\mu^{-1}c],
\]
where \(\theta_\mu\) is any element of multiplier \(\mu\).
The set \(H^1(C_2,H_G)\) has cardinality \(2^{\aleph_0}\).
For \(G=\widehat{GT}_0\), \(\Delta\) vanishes on this
closed family. All these \(H_G\)-conjugacy classes form
one \(G\)-conjugacy class.

\begin{proof}
A cocycle is \(a\in H_G\) with \(a\,c(a)=1\); it
corresponds to the involution \(ac\). Cohomology
identifies \(a\) and \(h^{-1}a\,c(h)\), exactly
conjugacy by \(H_G\). The Galois image makes
\(\lambda:G\to\widehat{\mathbb Z}^{\times}\) onto.
The element \(a_\mu=\theta_\mu c\theta_\mu^{-1}c\)
is a cocycle, and replacing \(\theta_\mu\) by a
different lift conjugates its involution by \(H_G\).

Two displayed involutions are \(H_G\)-conjugate
if and only if
\(\theta_\nu^{-1}h\theta_\mu\in C_G(c)=\{1,c\}\)
for some \(h\in H_G\), by A.6. Taking multipliers
gives \(\mu=\pm\nu\); these conditions are sufficient
as well. This proves the injection.
The cocycle space is compact, and its quotient by
the compact conjugating group is Hausdorff.
The orbit map from \(G\) factors continuously
through the multiplier quotient and then its
two signs, so the injection is continuous and closed.
Its domain has cardinality continuum. Conversely
\(G\) is contained in the metrizable profinite group
\(\operatorname{Out}(F)\), of cardinality at most
continuum. This proves equality.
Finally every represented involution is conjugate
to \(c\); the character \(\Delta\) is conjugacy
invariant and \(\Delta(c)=0\).
\end{proof}

\statement{A.8}{Proved}{Full conjugacy is an additional quotient of nonabelian cohomology}
With the notation of A.7, the full \(G\)-conjugacy
classes of nonidentity involutions are
\[
 \widehat{\mathbb Z}^{\times}\backslash H^1(C_2,H_G).
\]
The action of \(\mu\), using any lift \(\theta_\mu\), is
\[
 [a]\longmapsto[\,\theta_\mu a\,c(\theta_\mu^{-1})\,].
\]
For \(\mathcal N=\ker(\mathcal P\to P_F)\), the
\(\mathcal N\)-conjugacy classes of involutions with
image exactly \(c_F\) are \(H^1(C_2,\mathcal N)\).
Full \(\mathcal P\)-conjugacy on that fibre further
quotients by \(C_{P_F}(c_F)\), acting by lifted
conjugation, which need not preserve the distinguished
cohomology class.

\begin{proof}
By R1 every nonidentity involution has multiplier
\(-1\), so is uniquely \(ac\), \(a\in H_G\).
Its square-one condition is the cocycle equation.
Conjugation by \(\theta_\mu\) gives the formula.
Different lifts differ by \(H_G\), already divided
out in cohomology, proving the first assertion.

The same calculation for \(\mathcal P/\mathcal N\)
shows that involutions over the literal \(c_F\)
are \(ac\), \(a\in\mathcal N\). A conjugation
preserves this fibre exactly when its quotient
centralizes \(c_F\). Changing a lift changes
the action only by \(\mathcal N\)-conjugacy.
Such a lift need not centralize \(c\), explaining
the last qualification.
\end{proof}

\statement{A.9}{Proved}{Odd-prime conjugacy over a relative Frattini level}
Let \(U\) be open characteristic in \(F\), let \(p\)
be odd, put \(E_p=F/O^p(U)\), \(V=U/O^p(U)\),
and consider the peripheral outer subgroup preserving
\(V\). Its kernel on
\[
           E_p/\Phi(V)=F/\overline{U'U^p}
\]
is pro-\(p\). Two involutions with the same image
in this Frattini quotient are conjugate by the kernel.
This applies to the fibre over \(c\), where that
fibre is defined by descent to \(F/U\).

\begin{proof}
\(V\) is a nonabelian finitely generated free
pro-\(p\) group. Let \(J\) be the actual
automorphisms of \(E_p\) inducing identity on
\(E_p/\Phi(V)\). Restriction maps \(J\) into
the Frattini kernel of \(\operatorname{Aut}(V)\).
That kernel is pro-\(p\): on finite lower
\(p\)-central quotients its successive kernels
embed into homomorphism groups of elementary
abelian \(p\)-groups, and those quotients are cofinal.

Restriction is injective. If an element of \(J\)
fixes \(V\) pointwise, write \(a(g)=gu_g\) with
\(u_g\in\Phi(V)\). Comparing the conjugation
actions on \(V\) gives \(u_g\in C_V(V)=1\).
Thus \(J\) is pro-\(p\). If an outer automorphism
acts inner on \(E_p/\Phi(V)\), lift its inner
conjugator to \(E_p\) and correct it to an
element of \(J\). Consequently the outer
kernel is a continuous quotient of a closed
subgroup of \(J\), hence pro-\(p\).

The common image of the two involutions has
order two, since a pro-\(p\) group for odd \(p\)
has no involution. In its inverse image the
two cyclic subgroups are complements to the
normal pro-\(p\) kernel. Finite
Schur--Zassenhaus gives conjugators in every
finite quotient; compactness gives a conjugator
in the kernel itself.
\end{proof}

For arbitrary quotient automorphisms, preservation
of the image of a merely characteristic \(U\)
must be checked. It is included in the phrase
over the class of \(c\). For fully characteristic
\(U\) it is automatic: any quotient endomorphism
lifts to an endomorphism of the free profinite
group and preserves \(U\). Such \(U\) are cofinal.

\statement{A.10}{Proved}{Every negative pure pro-\(p\) involution is locally conjugate for odd \(p\)}
Let \(P^{(p)}\) be the pure peripheral outer group
of \(F^{(p)}\), with multiplier in
\(\mathbb Z_p^\times\), for odd \(p\).
Every involution of multiplier \(-1\) is
conjugate to \(c\) by a multiplier-one element.

\begin{proof}
The multiplier kernel \(H^{(p)}\) acts trivially
on abelianization, so lies in the pro-\(p\)
Frattini kernel used in A.9. The inverse image
of \(\{1,-1\}\) under the multiplier is
\(H^{(p)}\rtimes\langle c\rangle\).
Every negative involution defines a complement;
Schur--Zassenhaus gives the required conjugacy.
\end{proof}

This does not rescue the particular map of A.5.
Its local coboundary can be changed by a local
fixed element, whereas its full coboundary is
unique. The comparison of fixed subgroups
need not be surjective.

\statement{A.11}{Proved}{A peripheral reflection triangle is equivalent to full conjugacy}
For a nontrivial involution \(\alpha\in\mathcal P\)
let \(E=F\rtimes_\alpha C_2\). Then \(\alpha\)
is conjugate to \(c\) in \(\mathcal P\) exactly
when there are involutions \(t_0,t_1,t_2\)
representing its three section classes such that
\[
 a=t_1t_0,\qquad b=t_0t_2,\qquad d=t_2t_1
\]
generate conjugates of the respective inertia
groups of \(x,y,z\), and the \(t_i\) generate
\(E\) topologically. The last condition is
equivalent to generation in every \(E/R_U\).

\begin{proof}
For standard \(c\) take
\(t_0=c,t_1=xc,t_2=y^{-1}c\); their products
are \(x,y,z\) and they generate the extension.
A peripheral conjugator transports this triple.

Conversely the universal property gives a
surjection from the profinite coproduct
\(C_2^{(0)}\amalg C_2^{(1)}\amalg C_2^{(2)}\)
to \(E\), sending the factors to the named
involutions. Its parity kernel is free
profinite on \(a,b\). The restriction maps
onto \(F\): every even word in three
reflections is a word in the two edge products.
Both kernels are free profinite of rank two,
so Hopficity makes this restriction an
isomorphism. The equality \(abd=1\) is literal.
Abelianization makes the unit exponents in
the three peripheral conditions equal,
so the induced outer map belongs to \(\mathcal P\).
Conjugation by \(t_0\) inverts both \(a\) and \(b\).
Its outer action on \(F\) is \(\alpha\), since
\(t_0\) and the named splitting element have the
same image in \(E/F\). Hence the resulting outer
map conjugates \(c\) to \(\alpha\); no choice of
an actual splitting representative is needed.

Finally put \(H=\overline{\langle t_0,t_1,t_2\rangle}\).
If \(HR_U=E\) for all \(U\), then for every \(e\in E\)
the compact sets \(H\cap eR_U\) are nonempty
and have the finite intersection property.
Their intersection and \(\bigcap R_U=1\)
imply \(e\in H\). Thus \(H=E\).
\end{proof}

The transcription of cote 145, scan 23, lists
\(l_0=t_\infty t_1,\ l_1=t_0t_\infty,\
l_\infty=t_1t_0\). With ordinary multiplication
these satisfy \(l_0l_\infty l_1=1\).
The cyclic order or path-composition convention
must therefore be made explicit before
identifying them with \(x,y,z\).
We use \(abd=1\), and do not emend an unseen scan.

\statement{A.12}{Open}{Dyadic and mixed-prime involution obstructions}
A.9 does not determine all two-primary classes or all
gluing classes over \(U=F'F^3\) and \(U_{12}\).
A.25 below constructs an actual involution with
parity one in the image of \(\widehat{GT}_0\) at
\(U_{12}\), and A.24 proves that this parity
descends to that outer image. Thus the existence
assertion at this specified relative level is proved.
Lead 3 still needs a compatible order-two full lift.
Neither R2 nor A.3 supplies that further equivariance.

For a surjection \(G\to Q\) with kernel \(L\),
a chosen lift \(t\) of an involution can be
corrected to an involution \(ut\), \(u\in L\),
exactly when
\[
                    u(tut^{-1})=t^{-2}.
\]
Surjectivity asserts the existence of \(t\),
not solvability of this nonabelian norm equation.
The map \(C_4\to C_2\) disproves that implication
in general, but is not a counterexample in
the peripheral groups. The residual task is
to solve the norm equations compatibly in
the actual relative-level kernels with the
specified frames and \(\Delta=1\), or to
prove they are insoluble.

\statement{A.13}{Open}{Parity on arbitrary involutions}
For an involution \(\tau=(-1,f)\) in
\(\widehat{GT}_0\), R3 gives
\(\Delta(\tau)=\delta((1,f))\).
It is zero for A.7's continuum family and
for every \(\widehat{GT}_0\)-conjugate of \(c\).
The cocycle relation alone gives only
\(2\Delta(a)=0\) in \(\mathbb F_2\).
A.25 gives such an involution in an actual relative
outer image of \(\widehat{GT}_0\), but no full
involution of parity one or universal full vanishing
proof is obtained, in \(\widehat{GT}_0\) or
\(\widehat{GT}\).

\statement{A.14}{Open}{The three full conjugacy assertions}
We have not decided whether every involution
of \(\mathcal P,\widehat{GT}_0,\widehat{GT}\)
is conjugate to \(c\) in its respective group.
A.2 detects conjugacy of a fixed pair on
all relative levels; A.11 identifies the
equivalent peripheral generation condition.
Neither supplies the compatible local
conjugators or reflection triangles whose
existence is required.

For \S31, folios 198--199 verso, scans
391--394, the full finite-subgroup conclusion
proved by R1 is that the only abstract
possibilities in \(\mathcal P\) are \(1,C_2\),
and \(C_2=\langle c\rangle\) occurs.
Exhaustion of its embedding conjugacy
classes is exactly the remaining assertion.

\statement{A.15}{Proved}{The pronilpotent lifting kernel also has continuum many cohomology classes}
Put \(Q=\mathcal P_F\), \(\mathcal N=\ker(\mathcal P\to Q)\),
and write \(c_Q\) for the image of \(c\). There is a canonical
closed injection
\[
 C_Q(c_Q)/\langle c_Q\rangle
 \longrightarrow H^1(C_2,\mathcal N),\qquad
 \bar\beta\longmapsto[\beta c\beta^{-1}c].
\]
Its image is exactly the \(\mathcal N\)-conjugacy classes in the
fibre over \(c_Q\) which are conjugate to \(c\) in \(\mathcal P\).
Both its domain and \(H^1(C_2,\mathcal N)\) have cardinality
\(2^{\aleph_0}\). Thus the large cohomology set does not itself
provide an additional full involution-conjugacy class.

\begin{proof}
The pronilpotent quotient is the product of its maximal
pro-\(p\) quotients. Every factor is its characteristic
Sylow subgroup; conversely, a tuple of continuous
automorphisms of the factors gives a continuous automorphism
of their product. Inner automorphisms and the three peripheral
conditions are coordinatewise. Therefore
\[
 Q=\prod_p\mathcal P^{(p)},\qquad
 C_Q(c_Q)=\prod_p C_{\mathcal P^{(p)}}(c).
\]
For odd \(p\), the multiplier map on the latter centralizer
is onto \(\mathbb Z_p^\times\). Indeed choose an element
\(\theta\) of the desired multiplier using the Galois image.
By A.10, some multiplier-one \(v\) conjugates
\(\theta c\theta^{-1}\) to \(c\); then \(v\theta\)
centralizes \(c\) and retains that multiplier.
Consequently the displayed centralizer, and its quotient
by the finite group \(\langle c_Q\rangle\), have at least
continuum many elements.

By A.2 every \(\bar\beta\in C_Q(c_Q)\) lifts to
\(\beta\in\mathcal P\). Its displayed commutator with
\(c\) lies in \(\mathcal N\) and is a cocycle.
Changing the lift changes its class only by
\(\mathcal N\)-conjugacy. Two lifts \(\beta_1,\beta_2\)
give the same class exactly when
\(\beta_2^{-1}n\beta_1\in C_{\mathcal P}(c)=\{1,c\}\)
for some \(n\in\mathcal N\). On taking the quotient,
this is exactly equality of the two cosets modulo
\(\langle c_Q\rangle\); the converse follows by lifting
that equality. This proves injectivity and the asserted
description of the image.
Compactness makes the injection continuous and closed,
as in A.7. All groups and orbit spaces here are compact
metrizable, so their cardinalities are at most continuum.
The preceding lower bound proves the asserted equalities.
\end{proof}

If the same construction is made in \(\widehat{GT}_0\),
using its actual image in any \(c\)-stable quotient,
every class in the connecting-map image has \(\Delta=0\).
This follows from conjugacy invariance of the character;
no surjectivity onto the entire peripheral quotient is
required for this observation.

\statement{A.16}{Proved}{A prescribed reflection splits every pro-two extension of the symmetric group}
Let \(1\to K\to E\xrightarrow{q}S_3\to1\) be an extension
of profinite groups with \(K\) pro-two. If a transposition
\(t\in S_3\) has an involution lift \(\widetilde t\in E\),
then there is a continuous section \(j:S_3\to E\) satisfying
\(j(t)=\widetilde t\).

\begin{proof}
For every finite \(\mathbb F_2[S_3]\)-module \(M\), restriction
\[
 H^n(S_3,M)\longrightarrow H^n(\langle t\rangle,M)
\]
is an isomorphism for \(n\geq1\). Indeed the double-coset
formula for restriction after corestriction has the identity
term and one term through the trivial intersection
\(\langle t\rangle\cap r\langle t\rangle r^{-1}=1\),
where \(r\) is a three-cycle. The second term is zero
in positive degrees. Corestriction after restriction is
multiplication by the odd index three, hence the identity
on these \(\mathbb F_2\)-vector spaces.

First suppose \(K\) is finite. Induct on its order.
When \(K\ne1\), the subgroup
\(A=\{z\in Z(K):z^2=1\}\) is a nonzero elementary
abelian characteristic subgroup of \(K\), hence normal
in \(E\). Induction gives a section of \(E/A\to S_3\)
with the prescribed image of \(\widetilde t\).
Pull back that section to an extension with kernel \(A\).
Its class in \(H^2(S_3,A)\) restricts to zero, since
\(\widetilde t\) is an involution. The proved injectivity
of restriction makes the extension split.
The difference between any splitting on \(\langle t\rangle\)
and the specified lift is an \(A\)-valued one-cocycle.
Surjectivity in degree one extends its cohomology class to
\(S_3\); changing the section by that cocycle and then by
an \(A\)-conjugation makes its value at \(t\) literally
\(\widetilde t\). This proves the finite assertion.

For general pro-two \(K\), take all open normal subgroups
\(L\) of \(E\) contained in \(K\). Each finite quotient
has such a section with the prescribed transposition.
The corresponding closed subsets of the compact space
of tuples \((e_g)_{g\in S_3}\in E^{S_3}\), imposing the
section identities modulo \(L\), are nonempty and have
the finite intersection property. Their intersection
gives a section with all identities literal.
\end{proof}

\statement{A.17}{Proved}{The parity-one relative frame can be chosen anti-invariant under conjugation}
In the notation of R2, put \(H_0=F'F^3\),
\(\Gamma_2=\widehat B/O^2(H_0)\), and let
\(d=\operatorname{Inn}(s)\circ c\), so that
\(d(s)=s\), \(d(r)=r^{-1}\), \(d(x)=y^{-1}\),
and \(d(y)=x^{-1}\). There is an actual element
\[
 b\in C_{H_0^{(2)}}(r),\qquad q(b)=[5I],\qquad d(b)=b^{-1}.
\]
Consequently the identity automorphism of
\(\widehat B/R_{\mathrm{nil}}(H_0)\), with \(s\)- and
\(\tau\)-frames one and \(r\)-frame equal to \(b\) at two
and one at all other primes, is a framed cocycle for
conjugation by \(c\). Its modulo-twelve frame parity is one.
This establishes local equivariance of these data; it does
not assert an equivariant full lift.

\begin{proof}
Use R2's actual joint quotient
\[
 J=\operatorname{im}\bigl(\widehat B\to Q\times D_0\bigr),
 \qquad Q=\mathrm{SL}_2(\mathbb Z/12)/\{\pm I\},\quad
 D_0=\widehat B/H_0.
\]
The map \(\Gamma_2\to J\) has pro-two kernel; both factors
and the joint quotient are preserved by \(d\).
The element \(z_0=([5I],1)\in J\) is central of order two
and is fixed by \(d\). Thus
\(S=\langle\bar r,z_0d\rangle\) in \(J\rtimes\langle d\rangle\)
is an actual copy of \(S_3\), with its named three-cycle
and transposition.

There is a particularly small literal lift of that
transposition as an involution. Set
\[
 C=[x,y],\qquad h=CxCx^{-1}\in F'\subseteq H_0.
\]
In R2's signed matrix coordinates modulo twelve,
\[
 q(h)=\begin{pmatrix}9&4\\8&9\end{pmatrix},\qquad
 D=\begin{pmatrix}0&1\\1&0\end{pmatrix},\qquad
 q(h)\bigl(Dq(h)D^{-1}\bigr)^{-1}=5I.
\]
Conjugation by \(D\) represents \(d\) on \(Q\).
The \(D_0\)-coordinate of \(h\) is one, so the actual
involution \(hdh^{-1}\) maps to \(z_0d\).
Apply A.16 to the preimage of \(S\) in
\(\Gamma_2\rtimes\langle d\rangle\), retaining this
involution. Its lifted three-cycle is conjugate to
the named \(r\) by the pro-two kernel, by the
coprime complement-conjugacy argument of A.9.
After that conjugation write the lifted transposition
as \(bd\). Its image in \(D_0\rtimes\langle d\rangle\)
is \(d\), so \(b\in H_0^{(2)}\), and its \(Q\)-coordinate
gives \(q(b)=[5I]\). Its order two gives \(b\,d(b)=1\).
Since it inverts the named \(r\), while \(d\) also
inverts \(r\), the element \(b\) centralizes \(r\).

All components of the prescribed relative frame have
image one in \(D_0\), so they give an actual element
of the relative fibre product. For a multiplier-one
normalized automorphism \(\Phi\) of \(\widehat B\),
write \(\Phi(s)=\beta s\beta^{-1}\) and
\(\Phi(r)=A rA^{-1}\). Direct substitution from
\(c(s)=s\), \(c(r)=sr^{-1}s\) gives the frames of
\(c\Phi c\):
\[
 \beta^c=c(\beta),\qquad
 A^c=c(\beta)s\,c(\beta)^{-1}c(A)s.
\]
For the specified local data \(\beta=1\), the second
formula is \(A^c=d(A)=A^{-1}\). These are exactly
the frames of the inverse of the identity framed
automorphism. The parity is the same as in R2's E.58:
the signed \(r\)-frame is \(5I\), whose modulo-three
value is \(-I\), and the \(s\)-frame is one.
The independently executed certificate
\texttt{check\_A17\_frame.py} verifies the literal
matrix word and the covariance identities just used.
\end{proof}

The remaining full obstruction can now be placed in a
smaller, specified kernel. Let \(L\) consist of the
multiplier-one elements whose normalized action is
identity on \(\widehat B/R_{\mathrm{nil}}(H_0)\) and whose
two pure frames both have image one there. It is the
kernel of the framed action, is stable under \(c\), and
\(\Delta(L)=0\). R2 supplies a full \(\sigma\) lifting
A.17's data, with \(\Delta(\sigma)=1\). Put \(a=\sigma c\).
Then \(a^2\in L\). A parity-one involution in this
specified fibre exists exactly when some \(v\in L\)
solves
\[
                   v\,(ava^{-1})=a^{-2}.
\]
The local construction proves that the right-hand side
lies in this framed kernel. Solvability of this last full
norm equation remains open.

\statement{A.18}{Refuted}{R2 does not admit an unrestricted order-preserving refinement}
There is an actual involution of
\(\operatorname{Out}^{\sharp}(F/R_{\mathrm{nil}}(F))\)
which lifts to a full peripheral automorphism by R2 but
has no involution lift in \(\mathcal P\).

\begin{proof}
Use the actual product
\(F/R_{\mathrm{nil}}(F)=\prod_pF^{(p)}\).
Define \(\bar a\) to be standard conjugation \(c\)
on \(F^{(3)}\) and identity on every other factor.
This is a continuous automorphism of order two,
preserving each of the three peripheral conjugacy
classes up to the common unit
\[
 \lambda_3=-1,\qquad \lambda_p=1\quad(p\ne3).
\]
Its outer class is nonidentity, since its action on
the three-primary abelianization is \(-I\).
By A.2, which applies R2 after Galois translation
of the multiplier, this class has a lift in
\(\mathcal P\). Every lift has the displayed
multiplier, because the product quotient retains
the entire abelianization \(\widehat{\mathbb Z}^2\).
By R1, a nonidentity full involution must instead
have multiplier \(-1\) at every prime. Hence no
lift of \(\bar a\) can be an involution.
\end{proof}

This counterexample is in the actual peripheral level
groups, rather than in an unrelated finite extension.
It excludes a blanket equivariant strengthening of
the central-refinement count. A.17's remaining fibre
has multiplier globally \(-1\), so this orientation
obstruction does not decide that fibre.

\statement{A.19}{Proved}{The square obstruction at an abelian refinement and the character-lifting test}
Suppose \(a\) is a lift of a quotient involution,
\(a^2\in L\), and \(L_0\triangleleft\langle L,a\rangle\)
is contained in \(L\), with \(M=L/L_0\) finite abelian.
Conjugation by \(a\) makes \(M\) a \(C_2\)-module.
An involution exists above that quotient element in
\(\langle L,a\rangle/L_0\) exactly when
\[
 [a^2]=0\quad\hbox{in}\quad
 H^2(C_2,M)=M^{a}/(1+a)M.
\]
This obstruction vanishes when \(M\) has odd order.
For elementary abelian two-primary \(M\), its target is
\(\ker(a+1)/\operatorname{im}(a+1)\).

If \(\Delta\) has a continuous lift to a character
\(\widehat{GT}_0\to\mathbb Z/4\), then it vanishes on
every involution. The same conclusion follows from
a continuous lift of \(\delta\) on \(\mathcal H\)
which is invariant under conjugation by \(c\).
Existence of either lift is not established here.

\begin{proof}
The square of the conjugation action is conjugation
by \(a^2\in L\), which is trivial on \(M\).
Also \(a^2\) is invariant under conjugation by \(a\).
Replacing \(a\) by \(va\), with \(v\in L\), changes
its square in additive notation by \((1+a)\bar v\).
This proves the quotient criterion, which is the
usual cyclic two-cocycle description of \(H^2\).
For an invariant \(u\) of odd order, \(u/2\) is
defined and \((1+a)(u/2)=u\), proving vanishing.
In characteristic two the invariants are
\(\ker(a+1)\), giving the last formula.

A character to \(\mathbb Z/4\) sends an involution
to zero or two; both reduce to zero modulo two.
Alternatively, if \(b\in\mathcal H\) is a cocycle
for \(c\) and \(\theta\) is a \(c\)-invariant
lift of \(\delta\), then
\(0=\theta(bc(b))=2\theta(b)\), so
\(\delta(b)=0\).
\end{proof}

For the actual \(a=\sigma c\) of A.17, any proposed
global character lift \(\Theta\) would necessarily satisfy
\(\Theta(a^2)=2\), although \(a^2\in L\) and
\(\Delta(L)=0\). Thus such a character must detect the
residual square inside the framed kernel; it cannot be
inferred from the local frames alone. Conversely A.17's
local parity-one framed involution rules out a character
lift to \(\mathbb Z/4\) on that local framed group itself.

\statement{A.20}{Proved}{The locally glued symmetric subgroup does not lift in the standard extension}
Let \(H_0\), \(d\), and \(b\) be as in A.17. In the standard
relative extension
\[
 \bigl(\widehat B/R_{\mathrm{nil}}(H_0)\bigr)
                  \rtimes\langle d\rangle
\]
there is a subgroup \(S_{\mathrm{loc}}\simeq S_3\), generated
by the named \(r\) and a transposition with component \(bd\)
at two and \(d\) at every odd prime. No finite subgroup
of the standard full extension
\(\widehat B\rtimes\langle d\rangle\) maps onto
\(S_{\mathrm{loc}}\). The obstruction is visible already in
\(Q\rtimes\langle d\rangle\), with \(|Q|=576\).

\begin{proof}
All local components of the transposition have the same
image \(d\) in \(D_0\rtimes\langle d\rangle\).
They therefore form an actual element of the relative fibre
product. A.17 gives its order two and its inversion of \(r\).
Their images in \(D_0\rtimes\langle d\rangle\) generate
\(S_3\), so the resulting subgroup has order exactly six.

In the full free product \(\widehat B=C_2\amalg C_3\),
every element of order three is conjugate into the named
\(C_3\), and that factor is malnormal. These are
Herfort--Ribes Theorems A' and B'.\cite{A_HR1985}
In particular \(N_{\widehat B}(\langle r\rangle)=\langle r\rangle\).
Since \(d\) inverts \(r\),
\[
 N_{\widehat B\rtimes\langle d\rangle}(\langle r\rangle)
                         =\langle r,d\rangle\simeq S_3.
\]
Every full subgroup isomorphic to \(S_3\) intersects
\(\widehat B\) in its unique subgroup of order three:
it cannot be contained in \(\widehat B\), whose finite
subgroups conjugate into its factors by the same paper's
Theorem 1. Conjugating that order-three subgroup into
\(\langle r\rangle\) proves that every full \(S_3\)
is conjugate to \(\langle r,d\rangle\).

The image of \(S_{\mathrm{loc}}\) in the actual finite
quotient \(Q\rtimes\langle d\rangle\) is
\[
                   \langle R,[5I]d\rangle .
\]
It is not conjugate there to \(\langle R,d\rangle\).
Indeed a conjugator may be taken in \(Q\), since \(d\)
normalizes the latter subgroup, and it must normalize
its unique order-three subgroup. Directly solving the
two-by-two equations in \(\mathrm{SL}_2(\mathbb F_3)\)
shows that no determinant-one matrix conjugates \(R\)
to \(R^{-1}\), and its centralizer is \(\{\pm R^j\}\).
The centralizer modulo four is also \(\{\pm R^j\}\).
Traces exclude a negative sign in the projective
commutation equation. Thus a representative of a
normalizing element has coordinates
\(\varepsilon_4R^i\) modulo four and
\(\varepsilon_3R^j\) modulo three.
Conjugation by \(D\) fixes both signs and inverts \(R\).
Its \(d\)-norm consequently has positive signs in both
coordinates. It cannot lie in \([5I]\langle R\rangle\):
one global choice of representative has negative sign
modulo three, and the other has negative sign modulo four.
This proves the nonconjugacy and hence the asserted
nonlifting of an \(S_3\) subgroup.

For completeness, a finite subgroup mapping onto
\(S_{\mathrm{loc}}\) has trivial kernel: the kernel lies
in the torsion-free free profinite group \(H_0\).
It would therefore be an \(S_3\), already excluded.
The exact exhaustive certificate
\texttt{check\_A20\_standard\_extension.py} generates all
\(576\) elements of \(Q\), obtains normalizer and
centralizer orders \(18\), lists all nine \(d\)-norm
values and all three elements of the target coset, and
checks that the lists are disjoint. It independently
enumerates the two signed centralizers used above.
\end{proof}

A.20 blocks lifting this subgroup inside the \emph{standard}
extension. A hypothetical new involution of
\(\widehat{GT}_0\) would change the action defining the
full extension. No obstruction in that unknown twisted
extension is inferred from A.20.

\statement{A.21}{Refuted}{Real projectivity and a rank-two open free subgroup do not imply finite graph recognition}
There exists a finitely generated real-projective profinite
group \(E_*\) with an open normal subgroup isomorphic to
\(\widehat F_2\), quotient \(C_2\), and self-centralizing
involutions, which is not the profinite fundamental group
of a finite graph of finite groups. It has one conjugacy
class of involutions, so it does not satisfy R1's
three-section conclusion.

\begin{proof}
Put \(A=\prod_{p\text{ odd}}\mathbb Z_p\) and define
\[
 \Gamma_*=A\rtimes\mathbb Z_2,
 \qquad u\cdot a=(-1)^{u\bmod2}a.
\]
All groups and actions here are profinite and continuous.
The kernel of the parity map \(\Gamma_*\to C_2\) is
\(A\times2\mathbb Z_2\simeq\widehat{\mathbb Z}\).
The group \(\Gamma_*\) is torsion-free: a finite-order
element projects trivially to \(\mathbb Z_2\), and then
belongs to the torsion-free group \(A\).
Set
\[
 E_*=C_2\amalg\Gamma_*,\qquad
 \varepsilon:E_*\longrightarrow C_2,
\]
where \(\varepsilon\) is identity on the first factor
and parity on the second. In the open-subgroup Kurosh
decomposition of \(\ker\varepsilon\), the two vertex
groups are \(1\) and \(\widehat{\mathbb Z}\), and the
quotient graph has two vertices joined by two edges.
Consequently
\[
 \ker\varepsilon\simeq
 \widehat{\mathbb Z}\amalg\widehat{\mathbb Z}
 \simeq\widehat F_2.
\]
The open Kurosh theorem in precisely this full-profinite
form is recalled in Herfort--Ribes, Section~2.
Their Theorems A', B', and 1 show that all nonidentity
torsion lies in conjugates of the single \(C_2\)
factor, and that each such involution has centralizer
exactly that factor.\cite{A_HR1985}
Haran's characterization of real projectivity therefore
applies: an index-at-most-two subgroup of cohomological
dimension at most one, together with self-centralizing
involutions, suffices. We use its exact formulation in
P\'al--Quick, Theorem~2.1.\cite{A_PalQuickFormal2025}

Suppose \(E_*\) were the profinite fundamental group of
a finite graph of finite groups. Its vertex groups have
order at most two. Collapse each nonloop edge whose
nontrivial edge group is equal to an endpoint group.
In the resulting reduced graph, a loop with edge group
\(C_2\) would make its stable letter centralize that
\(C_2\); the stable letter has infinite order, as is
seen by killing the vertex groups and mapping that
loop to \(\widehat{\mathbb Z}\). This contradicts the
centralizer just proved. All remaining edges therefore
have trivial group. We obtain
\[
 E_*\simeq\bigl(\mathop{\amalg}_{j=1}^{k}C_2\bigr)
                   \amalg\widehat F_r .
\]
The single torsion class forces \(k=1\). The
torsion-free index-two subgroup of this free product
has rank \(2r+k-1\), by its two-sheeted graph or by
open Kurosh. The rank two already computed forces
\(r=1\).

On the other hand \([\Gamma_*,\Gamma_*]=A\), since
multiplication by two is invertible on \(A\).
Thus
\[
 E_*^{\mathrm{ab}}=C_2\times\mathbb Z_2,
\]
whereas \((C_2\amalg\widehat{\mathbb Z})^{\mathrm{ab}}
=C_2\times\widehat{\mathbb Z}\). The latter has
nontrivial odd-primary quotients and the former does not.
This contradiction proves the claim.
\end{proof}

The construction uses the torsion-free virtually
procyclic mechanism already present in Haran's
counterexamples to full-profinite virtual freeness
recognition.\cite{A_Haran1985}
It identifies a genuine limitation of that method;
it is not a new peripheral involution. For the groups
in A.14 the three-section hypothesis, and ultimately
the marked cusp products in A.11, still have to be used.

\statement{A.22}{Proved}{The three distinguished reflections normally generate, and a proper generated subgroup has infinite index}
Let \(\alpha\) and \(h\) be as in R1, and put
\(E=F\rtimes_\alpha\langle t\rangle\), \(b=hyh^{-1}\).
The three involutions
\[
       t_0=t,\qquad t_1=xt,\qquad t_2=b^{-1}t
\]
represent the three section classes and normally
generate \(E\). If their closed generated subgroup
\(J\) is proper, then \([E:J]\) is infinite.
Normal generation alone does not supply the full
generation required in A.11.

\begin{proof}
The equality \(f=h^{-1}\alpha(h)\) gives
\(\alpha(b)=b^{-1}\), so all three displayed
elements have order two. In
\(E^{\mathrm{ab}}\simeq C_2^3\), their images are
\(t,x+t,y+t\); these are distinct and hence the
elements are pairwise nonconjugate. R1 states that
there are exactly three classes. Quotienting \(E\)
by their closed normal closure kills \(t\), then
\(x=t_1t\), then \(b^{-1}=t_2t\). Since \(b\) is a
conjugate of \(y\), it kills \(y\), and the quotient
is trivial.

There is a continuous epimorphism
\(C_2\amalg C_2\amalg C_2\to J\) sending the three
factors to the displayed involutions. Its parity
kernel is free profinite of rank two, so \(J\cap F\)
is generated by at most two elements. If \(J\)
had finite index in \(E\), then \(J\cap F\) would
be open in \(F\) and the profinite Schreier formula
would give
\[
 d(J\cap F)=1+[F:J\cap F].
\]
It follows that \(J\cap F=F\), and \(t\in J\)
then gives \(J=E\). This proves the assertion.
\end{proof}

\statement{A.23}{Proved}{Torsion and cusp centralizers in the relative primary modular groups}
Let \(q:\widehat B=C_2\amalg C_3\to D\) be a finite
quotient injective on the two factors, with torsion-free
kernel \(U\), and write
\[
 \Gamma_p=\widehat B/O^p(U),\qquad V_p=U^{(p)}.
\]
For \(a=s\) or \(r\), of prime order \(\ell\), the
following statements hold:
\[
 \begin{array}{ll}
 p\ne\ell:&C_{\Gamma_p}(a)\longrightarrow C_D(q(a))
                    \text{ is onto},\\[2pt]
 p=\ell:&C_{\Gamma_p}(a)=\langle a\rangle.
 \end{array}
\]
For \(U=U_{12}=\ker q_{12}\) in the projective modular
quotient, put \(\Gamma=\widehat B/R_{\mathrm{nil}}(U)\).
Then
\[
 C_{\Gamma_p}(\tau)=\overline{\langle\tau\rangle},
 \qquad C_\Gamma(\tau)=\overline{\langle\tau\rangle}
                    \simeq\widehat{\mathbb Z}.
\]
In the second formula the closure is taken in \(\Gamma\);
in the first it is taken in \(\Gamma_p\).

\begin{proof}
If \(p\ne\ell\), choose a lift \(g\in\Gamma_p\)
of an element of \(C_D(q(a))\). The subgroups
\(\langle a\rangle\) and
\(\langle gag^{-1}\rangle\) are complements to
\(V_p\) in the inverse image of \(\langle q(a)\rangle\).
Coprime Schur--Zassenhaus, first in finite quotients
and then by compactness, conjugates the second to
the first by an element of \(V_p\). The generators
with image \(q(a)\) then agree. Adjusting \(g\)
by that element gives the asserted centralizer lift.

Now let \(p=\ell\), and put
\(J=q^{-1}(\langle q(a)\rangle)\subset\widehat B\).
The open Kurosh decomposition expresses \(J\) as
a finite free profinite product of copies of
\(C_p\), together with a finitely generated free
profinite factor. It has no factor of the other
prime order. Since \(J/U=C_p\),
\[
          J/O^p(U)=J^{(p)}.
\]
Indeed the left side is pro-\(p\), and every
homomorphism from \(J\) to a finite \(p\)-group
kills \(O^p(U)\). Its maximal pro-\(p\) quotient
is the corresponding free pro-\(p\) product.

If \(g\in C_{\Gamma_p}(a)\), choose
\(g_0\in\widehat B\) with the same \(D\)-image.
The elements \(a\) and \(g_0ag_0^{-1}\) belong
to \(J\) and become conjugate by \(V_p\) in
\(J^{(p)}\). Distinct Kurosh factors remain
nonconjugate in that free pro-\(p\) product,
by Herfort--Ribes Theorem B'. Their double-coset
labels in \(J\backslash\widehat B/\langle a\rangle\)
must consequently agree. Thus \(g_0\in J\),
so \(g\in J^{(p)}\). The same theorem gives
\(C_{J^{(p)}}(a)=\langle a\rangle\).
This proves the prime-primary assertion without
assuming a full-profinite Kurosh theorem for
arbitrary closed subgroups.\cite{A_HR1985}

For the cusp assertion, \(|D|=576\) and
\(q(\tau)\) has order twelve. The discrete kernel
has 48 cusps, genus 25, and free rank 97:
the cusp count is \(576/12\), and its Euler
characteristic is \(576(-1/6)=-96\).
In \(H_1(U,\mathbb F_p)\), the 48 oriented cusp
vectors have precisely the relation that their
sum is zero. In particular they are pairwise
distinct for every prime, including two and three.
The \(D\)-stabilizer of the vector represented
by \(\tau^{12}\) is exactly
\(\langle q(\tau)\rangle\): deck transformations
permute these vectors by their cusp permutation.
These assertions follow directly from the usual
oriented surface presentation of this torsion-free
finite cover of the modular orbifold.

If \(g\in\Gamma_p\) centralizes \(\tau\), it
centralizes \(\tau^{12}\in V_p\). Its image in
\(D\) therefore belongs to the indicated cusp
stabilizer. After multiplying \(g\) by an integral
power of \(\tau\), it belongs to \(V_p\).
The boundary element \(\tau^{12}\) is primitive
in \(U\), since another cusp can be omitted in
the surface free basis, and remains primitive
in \(V_p\). Herfort--Ribes Theorem B' in that
free pro-\(p\) group gives
\(C_{V_p}(\tau^{12})=\overline{\langle\tau^{12}\rangle}\).
This proves the first cusp equality. The fibre
product of these procyclic groups over
\(\langle q(\tau)\rangle=C_{12}\) is exactly
\(\widehat{\mathbb Z}\): its \(p\)-coordinate
is the infinite \(p\)-primary exponent, and the
remaining finite coordinates are the common
residue modulo twelve. This proves the second.
The same proof applies to \(x=\tau^{-2}\): its
negative sixth power is \(\tau^{12}\), so a
centralizer of \(x\) fixes the same cusp vector.
Conjugating by \(s\) treats \(y=sxs\). Thus
\(C_\Gamma(x)=\overline{\langle\tau\rangle}\)
and \(C_\Gamma(y)=s\overline{\langle\tau\rangle}s\),
as used below.
\end{proof}

\statement{A.24}{Proved}{The parity character descends to the actual outer image at level twelve}
Let
\[
 F_{12}=F/R_{\mathrm{nil}}(U_{12}),\qquad
 G_{12}=\operatorname{im}\!\left(
 \widehat{GT}_0\longrightarrow
               \operatorname{Out}^{\sharp}(F_{12})\right).
\]
The map is well-defined by the finite covariance in
refereed LM7\_PUSH E.65, which preserves \(U_{12}\).
There is a unique continuous character
\(\Delta_{12}:G_{12}\to\mathbb F_2\) such that
\(\Delta=\Delta_{12}\circ\pi_{12}\).

\begin{proof}
First consider a normalized multiplier-one action
\(\Phi\) which is identity on
\(\Gamma=\widehat B/R_{\mathrm{nil}}(U_{12})\).
Its pure \(s\)-frame, when projected to \(\Gamma_2\),
centralizes \(s\), so A.23 puts it in
\(\langle s\rangle\). Purity makes this projection
identity, since \(F\cap\langle s\rangle=1\)
already in the quotient \(S_3\). Hence its
\(Q_{12}\)-image is identity. At \(p=3\) the
same argument applies to the pure \(r\)-frame.
Moreover \(\Phi(s)=sf\) gives \(q_{12}(f)=1\).
The explicit covariance formula of LM7\_PUSH
E.65.9 then gives \(\rho_\Phi\equiv0\pmod3\),
and E.65.11 gives \(\rho_\Phi\in2\widehat{\mathbb Z}\).
Thus \(V_\Phi=I\) modulo twelve. Both parity
bits are zero. This proves that \(\Delta\)
vanishes on the kernel of the actual \(B\)-action.

We check that passing to an outer action on
\(F_{12}\) adds no kernel. Its abelianization is
\(\widehat{\mathbb Z}^2\), because
\(R_{\mathrm{nil}}(U_{12})\subset U_{12}'\subset F'\).
An outer-trivial action therefore has multiplier
one. Suppose its restriction is conjugation by
\(u\in F_{12}\). The centralizer of \(F_{12}\)
in \(\Gamma\) is trivial. To see this, the cusp
argument in A.23 also gives
\(C_\Gamma(x)=\overline{\langle\tau\rangle}\)
and
\(C_\Gamma(y)=s\overline{\langle\tau\rangle}s\).
Their images in \(S_3\) are distinct order-two
subgroups; an element of their intersection lies
in \(F_{12}\), and then in both
\(\overline{\langle x\rangle}\) and
\(\overline{\langle y\rangle}\). Abelianization
shows that their intersection is trivial.

An automorphism of \(\Gamma\) fixing the open
normal subgroup \(F_{12}\) pointwise is consequently
identity: for each \(g\in\Gamma\), its discrepancy
\(\Phi(g)g^{-1}\) centralizes that subgroup.
Thus the given \(\Phi\) is conjugation by \(u\)
on all of \(\Gamma\). It fixes \(\tau\), so A.23
gives \(u\in\overline{\langle\tau\rangle}\cap F_{12}
=\overline{\langle x\rangle}\). Write \(u=x^a\).
The normalized formula \(\Phi(s)=sf\) then gives
\(f=y^a x^{-a}\). Since \(f\in F_{12}'\), its
abelian coordinates \((-a,a)\) vanish; hence
\(a=0\). The actual action was identity after all.

We have proved kernel containment for \(\Delta\).
It therefore factors through \(G_{12}\); continuity
follows because a continuous surjection of compact
groups to a Hausdorff group is a quotient map.
\end{proof}

\statement{A.25}{Proved}{An actual parity-one involution exists at the relative level twelve}
The group \(G_{12}\) of A.24 contains an involution
\(a_{12}\) with
\[
             \lambda(a_{12})=-1,
             \qquad\Delta_{12}(a_{12})=1.
\]
It is not conjugate to the image of \(c\) in
\(G_{12}\). Thus the first, specified-relative-level
assertion of Lead~3 holds at \(U_{12}\).
The existence of an involution lifting it in the
full group is not asserted.

\begin{proof}
Let \(r_0=(3/4)^{1/4}>0\). The splitting field
\(L=\mathbb Q(i,r_0)\) has degree eight: the
polynomial \(4X^4-3\) is Eisenstein at three,
and \(\mathbb Q(r_0)\) is real. Choose a Galois
element \(\gamma\) whose restriction to \(L\)
fixes \(i\) and sends \(r_0\) to \(ir_0\).
Such a restriction is an automorphism of \(L\)
and extends to \(\overline{\mathbb Q}\).
Use its actual Galois image in \(\widehat{GT}\),
and put
\[
 \sigma_0=\gamma c\gamma^{-1}c,
 \qquad \sigma=\sigma_0^3.
\]
Then \(\lambda(\sigma_0)=1\) and
\(c\sigma_0c=\sigma_0^{-1}\), so
\(\sigma c\) is an actual full involution.
Also \(\sigma_0(r_0)=-r_0\). The Galois comparison
in refereed LM7\_PUSH E.65.13, including the model
and sign calculation supplied by its referee,
makes both normalized CM bits of \(\sigma_0\) equal
to one. More precisely these are
\(\varepsilon_3(T_{\sigma_0})\) and
\(\varepsilon_2(W_{\sigma_0})\), with the shear
\(V\) included in the definitions of \(T,W\).
LM7\_PUSH E.65.12 proves that \(T,W\) are
homomorphisms on \(\mathcal H\); composing with
the two indicated characters proves that both
bits are characters on all of \(\mathcal H\).
The same bits therefore hold for \(\sigma\).

The cubing calculation in the last paragraph of
that refereed proof gives
\[
 q\Phi_\sigma=q,
 \qquad q_{12}(f_\sigma)=I,
 \qquad V_\sigma=I.
\]
Let \(A,\beta\in F\) be its actual pure frames,
and write \(a=q(A),\ b=q(\beta)\). The r-bit
and s-bit are both one, while \(a\) and \(b\)
centralize \(q(r)\) and \(q(s)\), respectively.
More precisely, cubing the two finite cyclic
characters gives
\[
 a=[5I],\qquad
 b=[B]\ \text{or}\ [B^{-1}],\qquad
 B=\begin{pmatrix}9&4\\8&9\end{pmatrix}\pmod{12}.
\]
Here \(B\equiv I\pmod4\), \(B\equiv S\pmod3\),
and \(B^2=5I\). The choice between the two
possibilities will not matter.

On \(\Gamma_p=\widehat B/O^p(U_{12})\), define
\[
 \phi_2=\Phi_\sigma,\qquad
 \phi_p=\operatorname{id}\quad(p\ne2).
\]
All induce identity on the shared finite quotient
\(Q_{12}\). They therefore define an actual
continuous automorphism \(\phi\) of their fibre
product \(\Gamma\). Each fixes \(\tau\).
The automorphism \(\phi c\) has square one:
at two it is \(\Phi_\sigma c\), and at every
other prime it is \(c\).

We now verify full peripheral gluing with literal
conjugators, the essential point of the construction.
At two, A.23 supplies \(k_2\in C_{\Gamma_2}(r)\)
with \(q(k_2)=a^{-1}\). Replace the r-frame
\(A_2\) by \(A_2k_2\). Its \(Q_{12}\)-image
is one. At every other prime take the identity
r-frame. These give a single pure r-frame of
\(\phi\), with \(Q_{12}\)-image one.

At two retain the s-frame \(\beta_2\), with
\(Q_{12}\)-image \(b\). For every odd prime,
A.23 supplies \(k_p\in C_{\Gamma_p}(s)\)
with \(q(k_p)=b\). These are frames for the
identity automorphism at those primes. Their
common image \(b\in q(F)\) makes the resulting
global s-frame pure. The \(\tau\)-frame is one
at every prime. Thus \(\phi\) preserves all
three conjugacy classes \(s,r,\tau\), with
the stated common quotient frames.

Apply R2 with these literal conjugators. It gives
an automorphism of \(\widehat B\) fixing \(\tau\)
and lifting \(\phi\), with pure frames whose
\(Q_{12}\)-images are \((1,b)\) in the order
\((r,s)\). This lift belongs to \(\mathcal H\).
Indeed write it as \(s\mapsto sf\),
\(r\mapsto f^{-1}r\). Its reduction has
\(f_2=f_\sigma\) and \(f_p=1\) at odd primes,
so its image in every primary abelianization of
\(F\) is zero. Since
\(R_{\mathrm{nil}}(U_{12})\subset F'\), the
full \(f\) lies in \(F'\). The two torsion
relations are exactly the gentle relations, and
R2 supplies invertibility.

Denote this full multiplier-one lift by
\(\widetilde\sigma\). Its action on \(Q_{12}\)
is identity and its \(f\)-image there is one;
the covariance formula makes \(V=I\). Its
r-bit is zero and its s-bit is one. Hence
\(\Delta(\widetilde\sigma)=1\). The outer
image
\[
                  a_{12}=\pi_{12}(\widetilde\sigma c)
\]
is precisely the involution already constructed.
Its action on \(F_{12}^{\mathrm{ab}}\) is \(-I\),
so it has order exactly two. A.24 gives
\(\Delta_{12}(a_{12})=1\), whereas
\(\Delta_{12}(\pi_{12}(c))=0\). The character
separates their conjugacy classes.
\end{proof}

The exact finite computations underlying the construction
are retained in \path{check_A25_level.py} and its output:
the quotient has order 576, the pure subgroup has order 96,
the cusp orbit has size 48, and the frame matrix has
square \([5I]\). The script also verifies the dihedral
action on the four roots used above. The centralizer
and lifting assertions are proved in the text, not
inferred from this finite computation.

Unlike A.17, A.25 is an actual involution in an
outer image of the relative level. It uses the
different primes where torsion centralizers acquire
extra lifts, and retains their common finite quotient
frames throughout. Let
\(L_{12}=\ker(\widehat{GT}_0\to G_{12})\) and
\(a=\widetilde\sigma c\). Then
\(a^2\in L_{12}\) and \(\Delta(L_{12})=0\).
A full involution above this specified class exists
if and only if
\[
          v(ava^{-1})=a^{-2}
          \quad\text{has a solution }v\in L_{12}.
\]
A.19 describes every abelian refinement of this
remaining square obstruction. No solution or
incompatible finite refinement has yet been proved.

\statement{A.26}{Proved}{Finite frame visibility prevents arithmetic-only continuation}
Put
\[
 U=U_{12},\qquad V=\overline{U'U^6},\qquad D_6=\widehat B/V.
\]
This is a finite quotient of order \(576\cdot6^{97}\),
and \(V\) is characteristic in \(U\), normal in
\(\widehat B\), and preserved by \(\widehat{GT}_0\).
The two normalized CM bits of an element of
\(\mathcal H\) are determined by its actual action
on \(D_6\). In particular both bits are zero for
an element acting identically on \(D_6\).

Let \(H_6\) be the actual image of \(\mathcal H\)
in \(\operatorname{Aut}(D_6)\), and
\(J_6=\langle H_6,c_6\rangle\). There is a
character \(\Delta_6:J_6\to\mathbb F_2\) induced
by \(\Delta\) on
\(\mathcal H\rtimes\langle c\rangle\).
The involution of A.25 induces an actual involution
\(a_6\in J_6\) with \(\Delta_6(a_6)=1\).
No Galois element of multiplier \(-1\) has the
same actual action on \(D_6\).

\begin{proof}
For \(p=2,3\), write
\(D_p=\widehat B/\overline{U'U^p}\).
If \(a=s\) when \(p=2\), or \(a=r\) when
\(p=3\), then
\[
  \operatorname{im}\bigl(C_{D_p}(a)\to Q_{12}\bigr)
                      =\langle q(a)\rangle.
\]
To prove this, let \(g\) be a centralizing element
and choose a lift \(g_0\in\widehat B\) of its
\(Q_{12}\)-image. In
\(J=q^{-1}(\langle q(a)\rangle)\), the elements
\(a\) and \(g_0ag_0^{-1}\) become conjugate
by an element of \(U/\overline{U'U^p}\).
If their Kurosh factor labels were different,
the free-product universal property would give
a continuous character \(J\to C_p\) sending
one to one and the other to zero. Every such
character kills \(\overline{U'U^p}\), and
characters are constant on conjugacy classes.
This is a contradiction. Therefore the labels
coincide and \(q(g_0)\in\langle q(a)\rangle\),
as claimed. The reverse inclusion is immediate.

Any two pure s-frames for the same actual action
on \(D_6\) differ by a centralizer element.
The displayed assertion at two, followed by
intersection with the pure subgroup, shows that
their \(Q_{12}\)-images agree. The assertion at
three does the same for r-frames. The shear
\(V_\Phi\) modulo twelve is determined by the
\(Q_{12}\)-image of \(f\), by the explicit
covariance formula and the evenness of \(\rho\).
Thus both normalized bits are determined.

The kernel of \(\mathcal H\to H_6\) therefore
lies in \(\ker\delta\), giving a character
\(\delta_6\) on \(H_6\). It is invariant under
\(c_6\) because \(\delta\) is c-invariant.
The element \(c_6\) is not in \(H_6\):
on the retained abelian quotient
\(\widehat B^{\mathrm{ab}}=C_6\), its action
is inversion, whereas \(H_6\) acts identically.
Hence \(J_6=H_6\rtimes\langle c_6\rangle\),
and
\(\Delta_6(hc_6^e)=\delta_6(h)\) is well-defined.

The group \(U/R_{\mathrm{nil}}(U)\) has
characteristic subgroup whose quotient is
\(U/\overline{U'U^6}\). The actual involution
\(\phi c\) in A.25 therefore induces an actual
involution on \(D_6\). It belongs to \(J_6\),
has \(\Delta_6=1\), and is nonidentity since
it inverts the abelian quotient \(C_6\).
If a Galois element of multiplier \(-1\)
induced this action, it would belong to
\(\mathcal H\rtimes\langle c\rangle\) and
its \(\Delta\)-value would be one. This
contradicts R3's Galois vanishing.
\end{proof}

This is an obstruction to continuing the
\emph{arithmetic-only construction}, not an
obstruction to lifting \(a_{12}\) at all.
At any relative level below \(V\), the common
finite action must be \(a_6\). Thus no primary
component of such a continuation can be obtained
from a full Galois element of multiplier \(-1\).
In particular one cannot choose identity-on-
\(D_6\) Galois c-cocycles with the two CM bits
still equal to one. Nongalois relative lifts,
or a solution of the norm equation in
\(L_{12}\), remain possible. This conclusion
does not identify the entire character
\(\Delta\) on the image of arbitrary multipliers
in \(\operatorname{Aut}(D_6)\); the precise
finite domain here is \(J_6\).

\statement{A.27}{Proved}{The exact finite real-reflection count and its properness condition}
Let \(\pi:E\twoheadrightarrow D\) be an epimorphism of
finite groups. Fix involutions \(d_0,d_1,d_2\in D\)
which generate \(D\), and subsets
\(T_{ij}\subseteq\pi^{-1}(d_id_j)\), for
\((i,j)=(0,1),(1,2),(2,0)\). For every subgroup
\(J\leq E\), put
\[
 I_i(J)=\{u\in J:\pi(u)=d_i,\ u^2=1\},
 \qquad
 M_{ij}(J)_{u,v}=\mathbf1_{T_{ij}}(uv).
\]
Thus \(M_{ij}(J)\) is a rectangular matrix with rows
indexed by \(I_i(J)\) and columns indexed by
\(I_j(J)\). The number of triples of involutions
in \(E\) with the specified images and pair products,
and which generate \(E\), is
\[
 \sum_{J\leq E}\mu(J,E)
       \operatorname{tr}\bigl(M_{01}(J)M_{12}(J)M_{20}(J)\bigr),
 \tag{A.27.1}
\]
where \(\mu\) is the ordinary subgroup-lattice
M\"obius function. Empty index sets contribute zero.
This count concerns actual homomorphisms from
\(C_2*C_2*C_2\), with actual quotient data.

\begin{proof}
Expanding the trace gives exactly one summand for
\((u_0,u_1,u_2)\in I_0(J)\times I_1(J)\times I_2(J)\),
and that summand is one exactly when all three
pair-product conditions hold. Their ordered product
is automatically one:
\((u_0u_1)(u_1u_2)(u_2u_0)=1\).
Every counted triple has a unique generated subgroup.
M\"obius inversion on the finite subgroup lattice
therefore gives (A.27.1). The assumption that the
\(d_i\) generate makes every subgroup not mapping
onto \(D\) contribute zero.
\end{proof}

The matrices in (A.27.1) retain the common reflection
at the ends of adjacent edges. The ordinary
boundary-class trace sums of R2 do not retain this
incidence datum. Equality of those ordinary sums
therefore does not by itself establish equality of
(A.27.1). A concrete perfect-kernel control is
\(S_5\twoheadrightarrow C_2\), with kernel \(A_5\),
and all three \(d_i\) nonidentity. Every \(d_i\)
has an involution lift, and \(S_5\) has three
ordinary generators, but three such lifts cannot
generate \(S_5\): they are transpositions, whose
three edges cannot connect a graph on five vertices.
The certificate \path{check_A27_real_count.py}
checks perfectness of the kernel and enumerates
all 1000 triples. The generated-subgroup orders
and multiplicities are
\[
 (2,10),\quad(4,90),\quad(6,240),\quad
 (12,180),\quad(24,480).
\]
This is a control on a proposed general real
embedding theorem, not a counterexample to the
negative pants problem. Indeed a product of two
odd involutions in \(S_5\) has order at most three;
no generating pants triple of \(A_5\) has all
three orders at most three. The same certificate
checks all 816 candidate ordered pairs and finds
no pair generating \(A_5\).

\statement{A.28}{Proved}{A negative finite peripheral action is stopped by a central spin obstruction}
There is an epimorphism \(q:F\twoheadrightarrow A_6\)
and an involutory automorphism \(\vartheta\) of
\(A_6\) which sends each of the three labelled
cusp images to an \(A_6\)-conjugate of its inverse,
but no multiplier-minus-one peripheral automorphism
of \(F\) satisfies
\[
                       q\Phi=\vartheta q.
\]
In particular no full negative peripheral involution
induces this finite action. The obstruction is
already visible in a central double cover; it is
not a counterexample to R2 or A.3.

\begin{proof}
We use rightmost-first permutation composition and put
\[
\begin{aligned}
 X_0&=(1\,2\,3\,4)(5\,6),&
 Y_0&=(1\,2\,4\,6)(3\,5),\\
 Z_0&=(1\,2\,6\,3)(4\,5),&
 t&=(2\,4).
\end{aligned}
\]
The certificate \path{check_A28_s6_passport.py}
verifies \(X_0Y_0Z_0=1\) and
\(\langle X_0,Y_0\rangle=A_6\), giving
\(q(x)=X_0,q(y)=Y_0,q(z)=Z_0\).
Let \(\vartheta=\operatorname{Ad}(t)\).
The required inverse-class conjugators are,
in the displayed cusp order,
\[
                    1,\qquad(2\,6\,4),\qquad(2\,3\,4).
\]
In particular these conjugators are even.
Each cusp also has an odd involution which
inverts it: respectively \((2\,4),(2\,6),(2\,3)\).
The permutation \((3\,4\,6)\), of order three
in \(A_6\), cyclically permutes the three
cusp images by conjugation.

For completeness there is a short section-incidence
explanation of why an order-two full lift is impossible.
Odd involutions of \(S_6\) have just two cycle types,
\((2)\) and \((2,2,2)\). The product of two of
one type cannot have cycle type \((4,2)\).
For transpositions its order is at most three.
For two fixed-point-free involutions, their two
perfect matchings decompose into alternating cycles;
on six letters the possible product orders are
again at most three. Thus every \((4,2)\) cusp
requires opposite types at its two reflection ends.

In a hypothetical full negative pants extension,
R1 gives exactly three section classes. Their
images in
\(E^{\mathrm{ab}}=C_2^3\) are
\(\bar t,\bar x+\bar t,\bar y+\bar t\).
The x-cusp joins the first two, the y-cusp the
first and third, and the z-cusp the last two:
its two ends differ by \(\bar x+\bar y\).
Here each cusp has genuine involution ends.
Indeed an element of the odd coset which inverts
a full cusp has square in its procyclic centralizer;
inversion fixes that square, so the square is one.
Their three finite class labels would therefore
colour a triangle with two colours, with opposite
colours on each edge. This is impossible. No
assumption that the three reflections generate
is made in this argument.

A stronger obstruction excludes even infinite-order
full peripheral lifts of multiplier \(-1\).
We give a finite certificate rather than assume
that an arbitrary finite peripheral automorphism
lifts. The two degree-80 permutations in the
ATLAS representation \texttt{2S6G1-p80B0}
are opened source data. The code
\path{check_A28_double_cover.py} enumerates their
joint group with the natural standard generators
\((1\,2),(2\,3\,4\,5\,6)\) of \(S_6\).
It verifies an epimorphism
\[
 \widetilde S\twoheadrightarrow S_6,
 \qquad |\widetilde S|=1440,
 \qquad \ker=\langle\epsilon\rangle\cong C_2,
\]
with central kernel. Its even preimage
\(\widetilde A\) has order 720. Choose arbitrary
lifts \(X,Y\in\widetilde A\) of \(X_0,Y_0\),
and put \(Z=(XY)^{-1}\). Choose an involution
\(T\in\widetilde S\) above \(t\).
The enumeration of all conjugators in
\(\widetilde A\) proves, for each
\(Q\in\{X,Y,Z\}\),
\[
 TQT^{-1}\not\sim_{\widetilde A}Q^{-1},
 \qquad
 TQT^{-1}\sim_{\widetilde A}\epsilon Q^{-1}.
 \tag{A.28.1}
\]
Each lift has order eight. More precisely the
unit exponents \(k\pmod8\) for which
\(TQT^{-1}\sim Q^k\) are exactly \(3,5\).
The same statements hold for either central
choice of each lift. All 1440 elements, the
central kernel, the 720 even conjugators and
the stated class tests are checked by the code;
the group order is not inferred solely from
the ATLAS label.\cite{AAtlasA6}

Let \(q_2:F\to\widetilde A\) send \(x,y\)
to \(X,Y\). If \(\Phi\) as in the statement
existed, the two homomorphisms
\(q_2\Phi\) and \(\operatorname{Ad}(T)q_2\)
would have the same projection to \(A_6\).
Their ratio is consequently a continuous central
character \(\eta:F\to\langle\epsilon\rangle\).
The multiplier-minus-one cusp conditions and
(A.28.1) require
\[
             \eta(x)=\eta(y)=\eta(z)=\epsilon.
\]
But \(xyz=1\), contradicting
\(\eta(x)\eta(y)\eta(z)=1\).
This proves the stronger assertion.
\end{proof}

The kernel of \(q\) is open normal; it is not
asserted to be characteristic. Any characteristic
refinement retaining this quotient and this action
inherits the displayed obstruction. Conversely,
we have not constructed an automorphism of such a
characteristic refinement or a relative pronilpotent
level inducing \(\vartheta\). The central calculation
explains why the finite inverse-class data alone
cannot be inserted into R2. The dyadic class
obstruction, rather than perfectness alone, is
what is retained here.

\statement{A.29}{Proved}{The new relative involution has exactly three section classes}
For the involution \(a_{12}\) constructed in A.25,
choose its actual representative \(\phi c\) on
\(F_{12}=F/R_{\mathrm{nil}}(U_{12})\).
Then
\[
                \bigl|H^1(C_2,F_{12})\bigr|=3,
\]
where the action is this specified representative.
Its relative dyadic component likewise has exactly
three section classes and has trivial fixed subgroup
in its orientation kernel. These are section
classes of the group extension, not conjugacy
classes of peripheral automorphisms.

\begin{proof}
First let
\(W=C_2\amalg C_2\amalg C_2=F\rtimes\langle c\rangle\),
and let \(U\) be any open normal subgroup of
\(W\) contained in \(F\). In
\(W/O^2(U)\), every nonidentity torsion element
of order two has nonidentity image in \(W/U\),
because \(U^{(2)}\) is torsion-free. If that
image generates \(C\cong C_2\), let
\(J\leq W\) be its inverse image. The open
Kurosh theorem expresses \(J\) as a free
profinite product of order-two factors and a
free profinite group. Since \(J/U=C_2\),
\[
                     J/O^2(U)=J^{(2)}.
\]
Indeed the left side is pro-2, and every finite
2-group quotient of \(J\) kills \(O^2(U)\).
The torsion theorem for the resulting free pro-2
product shows that the given involution is conjugate
to an image of a full factor of \(W\).
The three full factors remain nonconjugate,
because \(W^{\mathrm{ab}}=C_2^3\) survives this
quotient. Conjugacy under the even subgroup
is the same as full conjugacy for odd involutions:
one may multiply an odd conjugator by its
source involution to make it even.

The same open-product argument proves that every
such torsion centralizer is exactly the cyclic
subgroup it generates. Explicitly, if a centralizing
element has finite image \(g\), conjugating its
factor by a lift of \(g\) gives two factors of
\(J\) that become conjugate in \(J^{(2)}\).
The free pro-2 torsion theorem forces equality
of their Kurosh labels, so \(g\in C\).
The centralizer computation now takes place inside
\(J^{(2)}\), where the free-factor centralizer
is that factor. Its intersection with the even
subgroup is trivial. This uses the free-factor
case of Herfort--Ribes, not a closed-subgroup
Kurosh assertion.\cite{A_HR1985}

Apply this to \(U=U_{12}\), which is preserved
by the actual Galois action used in A.25.
The dyadic involution \(\Phi_\sigma c\) is
Galois-conjugate to \(c\), so its section set
has the same three classes and trivial even
centralizers. The choice of conjugating Galois
element preserves \(U_{12}\), by R3's quotient
covariance.

At each odd prime the kernel above the common
finite quotient is pro-\(p\). Every involution
in the specified quotient coset has an involution
lift, and any two such lifts with the same finite
image are conjugate by that kernel, by
Schur--Zassenhaus. Moreover the centralizer of
an involution surjects onto the centralizer of
its finite image: conjugate a lift by any chosen
finite centralizer lift, then apply the same
complement-conjugacy assertion in the kernel.

It follows that projection to the dyadic component
induces a bijection on section-conjugacy classes
for the complete fibre product. Surjectivity
follows by choosing odd-prime lifts with the
same finite image. For injectivity, choose a
dyadic conjugator and its finite image; at every
odd prime use centralizer surjectivity to adjust
a conjugator to this very same finite image.
The resulting tuple is an element of the fibre
product and conjugates the two sections. This
proves the asserted cardinality three.
\end{proof}

The unmarked pro-2 recognition theorem is fully
consistent with A.29. Herfort--Zalesskii's
Theorem~6.2 conjugates a finite 2-subgroup of
\(\operatorname{Aut}(F^{(2)})\) into the
automorphisms of a dense abstract free group.
It does not retain three prescribed peripheral
classes.\cite{AHerfortZalesskii2013}
Thus neither that theorem nor the section count
proves the peripheral conjugacy statement left
open in A.12. The actual involution of A.25
already passes the three-section test.

\statement{A.30}{Proved}{The level-twelve frames can also satisfy the involution relations}
The construction in A.25 can be made with its literal
relative frames \(A,\beta\in\Gamma\) so that the
associated r-normalizer and s-normalizer in
\(\Gamma\rtimes\langle\phi c\rangle\) are actual
subgroups \(S_3\) and \(C_2\times C_2\), with the
finite frame values still
\[
                         q(A)=1,\qquad q(\beta)=b.
\]
Consequently an R2 lift \(a=\widetilde\sigma c\)
can be chosen with
\[
 a^2\in L_{12}^{\mathrm{fr}},\qquad
 L_{12}^{\mathrm{fr}}
 =\{k\in\mathcal H:\Phi_k|_\Gamma=1,
                  \ A_k|_\Gamma=\beta_k|_\Gamma=1\}.
 \tag{A.30.1}
\]
This is a smaller kernel than the unframed
\(L_{12}\). It does not assert a lift at a
smaller characteristic level or solve the remaining
norm equation.

\begin{proof}
Write \(a_2=\Phi_\sigma c\) for the genuine
Galois involution used at the dyadic component,
and write \(A_2,\beta_2\) for its underlying
multiplier-one element's frames. Put
\[
 J_2=\Phi_\sigma(s)A_2,
 \qquad D_2=J_2^{-1}a_2,
 \qquad T_2=\beta_2^{-1}a_2.
\]
The element \(D_2\) inverts \(r\), because
\(c(r)=sr^{-1}s\). In the full Galois extension,
its square lies in \(C_{\widehat B}(r)=\langle r\rangle\)
and is fixed by inversion; hence \(D_2^2=1\).
Similarly \(T_2\) centralizes \(s\) and its
square lies in \(\langle s\rangle\). That square
cannot be \(s\): the abelianization of
\(\widehat B\rtimes\langle\Phi_\sigma c\rangle\)
is \(C_2^2\), and \(s\) has nonzero image.
Thus \(T_2^2=1\). These equations descend to
\(\Gamma_2\rtimes\langle a_2\rangle\).

In the common finite extended quotient, put
\(d=sc\) and \(z_0=[5I]\). Since
\(q\Phi_\sigma=q\), the frame values give
\[
                       q(D_2)=z_0d,
                       \qquad q(T_2)=b^{-1}c.
\]
Choose a pure \(h\in\Gamma_2\) with finite
image \([B]\), where \(B\) is the matrix of
A.25. Direct multiplication gives
\[
             d(B)=B^{-1},\qquad B^2=5I.
\]
It follows that the involution \(hD_2h^{-1}\)
has finite image \(d\). The inverse image of
\(\langle q(r),d\rangle\cong S_3\) in this
relative extension has pro-2 kernel. A.16 supplies
a section retaining this exact involution. Its
order-three subgroup is conjugate to \(\langle r\rangle\)
by an element of \(U_{12}^{(2)}\), by coprime
complement conjugacy. After this conjugation the
section has an involution \(D'_2\) satisfying
\[
        D'_2r(D'_2)^{-1}=r^{-1},\qquad q(D'_2)=d.
\]
Then \(k_2=D_2D'_2\) centralizes \(r\), belongs
to \(\Gamma_2\), and has finite image \(z_0\).
Replace \(A_2\) by \(A'_2=A_2k_2\).
The r-frame now has finite image one, and its
associated r-normalizing reflection is \(D'_2\).
At every odd prime retain the identity r-frame
and the standard r-normalizing reflection \(sc\).
Thus the r-frames glue, and their normalizers
glue to an actual finite \(S_3\).

At two retain \(\beta_2\) and \(T_2\).
For an odd prime \(p\), A.23 gives a surjection
\[
             C_{\Gamma_p}(s)\longrightarrow C_{Q_{12}}(q(s))
\]
with pro-\(p\) kernel, equivariant for \(c\).
The element \(b^{-1}c\) has order two in the
finite semidirect product, since \(c(b)=b^{-1}\).
Schur--Zassenhaus supplies an involution above it
in \(C_{\Gamma_p}(s)\rtimes\langle c\rangle\).
Write this involution as \(\beta_p^{-1}c\).
Then \(q(\beta_p)=b\), and \(\beta_p\) is a
frame for the identity action at that prime.
These frames glue with \(\beta_2\); their
involutions glue to \(T=\beta^{-1}(\phi c)\),
which centralizes \(s\). Since \(T\) is outside
\(\Gamma\), it and \(s\) generate \(C_2^2\).
The unchanged \(\tau\)-frame is one.

Apply R2 to these actual frames, and call the full
lift \(a=\widetilde\sigma c\). It still has
\(\Delta(a)=1\) and the same relative action as
A.25. For its s-frame, applying \(a\) twice
gives the pure conjugator \(a(\beta)\beta\).
Its relative image is one, exactly because
\((\beta^{-1}(\phi c))^2=1\).
For the r-frame put \(J=\Phi_{\widetilde\sigma}(s)A\).
Applying \(a\) twice gives the pure conjugator
\(a(J)J\), whose relative image is one because
\((J^{-1}(\phi c))^2=1\). Uniqueness of the full
pure frames identifies these with the frames of
\(a^2\). Together with its identity relative
action, this proves (A.30.1).
\end{proof}

The kernel \(L_{12}^{\mathrm{fr}}\) is invariant
under conjugation by this chosen \(a\). One may
verify this by composing actual peripheral actions
and their literal conjugators: an identity action
with both conjugators equal to one remains such
under framed conjugation. In the specified
\emph{framed} fibre, an order-two full lift exists
exactly when
\[
             v(ava^{-1})=a^{-2}
             \quad\text{for some }v\in L_{12}^{\mathrm{fr}}.
 \tag{A.30.2}
\]
Solving (A.30.2) would solve Lead~3. Failure to
solve it would exclude only this framed fibre;
a full involution over the unframed \(a_{12}\)
could still use different relative frames, and
is governed by A.25's equation in \(L_{12}\).
Neither equation is declared solvable or insoluble.

\statement{A.31}{Proved}{Parity-one involutions are compatible through the entire relative principal-congruence tower}
For every positive integer \(N\) divisible by twelve, put
\[
 \begin{gathered}
 Q_N=\operatorname{SL}_2(\mathbb Z/N)/\{\pm I\},\qquad
 U_N=\ker(\widehat B\longrightarrow Q_N),\\
 \Gamma_N=\widehat B/R_{\mathrm{nil}}(U_N),\qquad
 F_N=F/R_{\mathrm{nil}}(U_N).
 \end{gathered}
\]
These kernels are preserved by the actual normalized
\(\widehat{GT}_0\)-action. Let \(G_N\) be its actual outer
image on \(F_N\). There is a single full element
\(a\in\widehat{GT}_0\) such that
\[
 \lambda(a)=-1,\qquad \Delta(a)=1,\qquad
 a_N^2=1\quad\hbox{for every }12\mid N,
 \qquad a_N=\pi_N(a)\in G_N.
 \tag{A.31.1}
\]
The \(a_N\) are compatible and each has order exactly two.
The character \(\Delta\) descends to every \(G_N\), and
\(\Delta_N(a_N)=1\). In particular their inverse limit
is an involution of parity one in the image of the full
group on this entire tower.

The assertion is restricted to this tower. The full
element \(a\) is not proved to have order two, and the
relative principal-congruence tower is not cofinal in
the tower in A.2.

\begin{proof}
We first choose the arithmetic cocycle more carefully
than in A.25. Fix an embedding
\(\overline{\mathbb Q}\hookrightarrow\overline{\mathbb Q}_3\)
and write \(r_0=(3/4)^{1/4}>0\). The field
\(\mathbb Q_3(i)\) is unramified quadratic over
\(\mathbb Q_3\); over it the polynomial
\(4X^4-3\) is Eisenstein. Since it contains the fourth
roots of unity, adjoining \(r_0\) gives a cyclic,
totally and tamely ramified extension of degree four.
Its inertia group therefore contains the automorphism
\(r_0\mapsto ir_0\), fixing \(i\). Choose a lift
\(\gamma\) in the absolute inertia subgroup at three.

Every splitting field obtained by adjoining roots of
unity of 2-power order and a 2-power root of two is
unramified at three: the discriminant of
\(X^{2^m}-2\) is a power of two, and the same holds
for the relevant cyclotomic extensions. Thus
\(\gamma\) fixes all these roots. Let
\[
 \sigma_0=\gamma c\gamma^{-1}c.
\]
It has multiplier one, satisfies
\(c\sigma_0c=\sigma_0^{-1}\), sends \(r_0\) to
\(-r_0\), and has zero dyadic Kummer character of two.
The last statement also follows directly because both
\(\gamma\) and \(c\) fix each positive real
\(2^{1/2^m}\).

Let \(e_2\in\widehat{\mathbb Z}\) be the idempotent
whose 2-adic component is one and whose odd-primary
components are zero. Define, in the procyclic closure
of \(\sigma_0\),
\[
                         \sigma=\sigma_0^{e_2}.
\]
The complete Kummer character \(\kappa_2\) is additive
on the multiplier-one Galois subgroup. Hence
\[
 \kappa_2(\sigma)=e_2\kappa_2(\sigma_0)=0,
 \qquad c\sigma c=\sigma^{-1},
 \qquad \sigma(r_0)=-r_0.
 \tag{A.31.2}
\]
The Galois comparison in R3 again makes both normalized
CM bits equal to one. Moreover \(\sigma c\) is a
Galois conjugate of \(c\): if
\(b=(e_2-1)/2\in\widehat{\mathbb Z}\), then
\[
 \sigma c=(\sigma_0^b\gamma)c(\sigma_0^b\gamma)^{-1}.
\]

Marin's modular representation calculation identifies
its congruence shear on a Galois element with
\(-8\kappa_2\), prime by prime
\cite[\S4, p.~776]{AMarin2008}.
Consequently
\[
                         q_N\Phi_\sigma=q_N
                         \quad\text{for every }N.
 \tag{A.31.3}
\]
Here the sign or left-versus-right normalization of
the shear is irrelevant to its vanishing. The full
covariance formula in refereed E.65.15 proves that
every \(U_N\) is preserved by
\(\widehat{GT}_0\). We do not assert that \(U_N\)
is characteristic under all automorphisms of \(F\).

Write \(A_\sigma,\beta_\sigma\) for the full pure
r- and s-frames of \(\sigma\), and put
\(a_N'=q_N(A_\sigma)\), \(b_N=q_N(\beta_\sigma)\).
At twelve, the signed r-character takes values in
\(C_6\), and the signed s-character in \(C_4\).
Since \(e_2\equiv3\pmod6\) and
\(e_2\equiv1\pmod4\), (A.31.2)--(A.31.3) give
\[
 a_{12}'=[5I],\qquad b_{12}=[B]^{\pm1},\qquad
 B=\begin{pmatrix}9&4\\8&9\end{pmatrix}\pmod{12}.
 \tag{A.31.4}
\]

For fixed \(N\), use the fibre-product decomposition
\[
 \Gamma_N=\mathop{\prod}_{Q_N,p}
          \Gamma_{N,p},\qquad
 \Gamma_{N,p}=\widehat B/O^p(U_N).
\]
Define the actual automorphism \(\phi_N\) by
\(\Phi_\sigma\) on its dyadic component and identity
on every odd-primary component. Equation (A.31.3)
gives the common identity action on \(Q_N\), so
these actions glue; \((\phi_Nc)^2=1\).
They are compatible as \(N\) varies, since the same
full \(\sigma\) is used at each level.

The literal frames glue as well. At two, the
coprime-centralizer assertion of A.23 gives
\(k_2\in C_{\Gamma_{N,2}}(r)\) of finite image
\((a_N')^{-1}\). Replace the dyadic r-frame by
\(A_{\sigma,2}k_2\); its finite image is one.
Use the identity r-frame at every odd prime.
For s, retain \(\beta_{\sigma,2}\) at two.
At every odd prime A.23 supplies an s-centralizing
frame with finite image \(b_N\). All these frames
are pure because their common finite images are
in \(q_N(F)\) and \(U_N\subset F\).
The \(\tau\)-frame is one. R2 therefore supplies
a full lift in \(\mathcal H\), just as in A.25.
The commutator normalization follows because its
\(f\)-image has zero abelianization at every prime
and \(R_{\mathrm{nil}}(U_N)\subset F'\).

For compatibility of the full lifts, define the closed
subsets of the compact group \(\mathcal H\)
\[
 \mathcal C_N=\{h:\ \Phi_h|_{\Gamma_N}=\phi_N,
       \ q_N(A_h)=1,\ q_N(\beta_h)=b_N\}.
\]
The preceding construction proves each nonempty.
For any finite set of moduli, a common multiple
\(M\) satisfies \(\mathcal C_M\subset\bigcap_N
\mathcal C_N\). Compactness gives
\(h\in\bigcap_{12\mid N}\mathcal C_N\).
This argument requires no surjectivity assertion
for transition maps between centralizers.
Set \(a=hc\). Its action on every \(\Gamma_N\)
is \(\phi_Nc\), of square one. The frame values
at twelve in (A.31.4), now changed to \((1,b_{12})\),
give \(\Delta(h)=1\); also \(\Delta(c)=0\).
Every \(a_N\) acts by \(-I\) on
\(F_N^{\mathrm{ab}}=\widehat{\mathbb Z}^{,2}\),
so it has order exactly two. Since \(U_N\subset U_{12}\),
the outer image \(G_N\) maps to \(G_{12}\).
A.24 proves that \(\Delta_N\) is well defined.
This proves (A.31.1).

We finally verify the precise limitation on the tower.
Every finite quotient of \(\Gamma_N\) is an extension
of a quotient of \(Q_N\) by a finite nilpotent group.
Its nonabelian composition factors are therefore of
the form \(\operatorname{PSL}_2(\mathbb F_p)\),
with \(p\) a prime. The same restriction holds for
finite quotients of the open normal subgroup \(F_N\):
given such a quotient, take the finite intersection
of the conjugates of its kernel in \(\Gamma_N\),
then use the resulting finite quotient of
\(\Gamma_N\). A normal subgroup and its quotient
cannot acquire new composition factors.

There is an epimorphism \(F\twoheadrightarrow A_7\),
sending \(x\) to \((1\,2\,3\,4\,5\,6\,7)\)
and \(y\) to \((1\,2\,3)\). These generate
\(A_7\), which is simple of order 2520.
It is not \(\operatorname{PSL}_2(\mathbb F_p)\):
for \(p\ge5\), equality of orders would require
\(p\mid2520\), leaving \(p=5,7\), whose respective
orders are 60 and 168. The cases \(p=2,3\)
are not nonabelian simple. The finite certificate
\path{check_A31_congruence_scope.py} verifies
generation, simplicity by all conjugacy-class normal
closures, and these order calculations.

Put \(T=\bigcap_{12\mid N}R_{\mathrm{nil}}(U_N)\).
Compactness gives
\(F/T\cong\varprojlim_NF_N\), and every finite
quotient of this inverse limit factors through one
\(F_N\). Thus \(A_7\) is not its quotient.
The image of \(T\) under the displayed epimorphism
to \(A_7\) is consequently nontrivial, and hence
is all of \(A_7\). In particular \(T\ne1\),
proving noncofinality. This group-theoretic kernel
statement alone makes no assertion about the
faithfulness of the \(\widehat{GT}_0\)-action
on the quotient.
\end{proof}

The sharper arithmetic-only obstruction of A.26 is
consistent with this construction. Its finite
quotient \(D_6\) is not a quotient of any of the
common bases \(Q_N\) through the named quotient
maps: the \(\sigma\) above acts identically on all
\(Q_N\), but its two nonzero normalized CM bits
prevent identity on \(D_6\) by A.26.
Thus no \(U_N\) is contained in
\(\overline{U_{12}'U_{12}^6}\).
This concerns the finite common bases; \(D_6\)
is already a quotient of the larger relative group
\(\Gamma_{12}\). A continuation below that kernel
would require additional information beyond A.31.

\statement{A.32}{Proved}{Congruence CM refinements cannot supply the invariant fourth-root character obstruction}
Let \(\mathcal H_{\mathrm{cong}}\) be the image of
\(\mathcal H\) on the family of actual groups
\(\Gamma_N\), for \(12\mid N\). The character
\(\delta=\Delta|_{\mathcal H}\) descends to this
image. There is no continuous homomorphism
\[
 \eta:\mathcal H_{\mathrm{cong}}\longrightarrow
             \mathbb Z/4
\]
whose reduction modulo two is \(\delta\) and
which is invariant under conjugation by \(c\).
The same exclusion applies to a proposed character
on \(\mathcal H\) depending only on its congruence
CM-frame data, including any refinement at moduli
24 or 48, or continuously on their entire tower.

\begin{proof}
The descent of \(\delta\) follows from A.24.
Let \(h\) be the element constructed in A.31,
and let \(\bar h\) be its actual image on the
entire relative congruence tower. There
\(\bar h(c\bar h c)=1\), while
\(\delta(\bar h)=1\). A putative \(\eta\)
would give
\[
 0=\eta\bigl(\bar h(c\bar h c)\bigr)
   =2\eta(\bar h)=2\quad\text{in }\mathbb Z/4,
\]
a contradiction.

For the frame formulation, two normalized full
elements having the same actual action on
\(\Gamma_N\) have r-frames differing by an
r-centralizer there. Projecting to
\(\Gamma_{N,3}\), A.23 identifies that centralizer
with \(\langle r\rangle\); purity makes the
difference trivial. Their finite r-frame images
therefore agree. The same argument at two works
for the s-frames. The finite signed frame is
determined by its projective class, since its
reduction modulo four lies in the pure four-group,
which does not contain \(-I\).
The multiplier-one congruence shear is likewise
determined by the actual \(Q_N\)-action: a shear
\(U^t\) centralizing the projective image of
\(S\) has \(t=0\pmod N\), by its lower-left
and diagonal entries. Hence the normalized CM
values formed from these frames and the shear
are functions of the actual action.
Any character depending only on those values
factors through \(\mathcal H_{\mathrm{cong}}\).
The compact quotient topology makes this factor
continuous, so the preceding contradiction applies.
\end{proof}

A global \(c\)-invariant \(\mathbb Z/4\)-lift of
\(\delta\), if it exists, must therefore use data
not determined by this relative congruence tower.
Neither its existence nor a solution of the full
norm equation is proved here.

\statement{A.33}{Proved}{The compatible congruence family also has equivariant literal frames}
The full element \(a=hc\) in A.31 can be chosen so that
its square belongs to
\[
 L_{\mathrm{cong}}^{\mathrm{fr}}=
 \bigcap_{12\mid N}\{k\in\mathcal H:
       \Phi_k|_{\Gamma_N}=1,\quad
       A_k|_{\Gamma_N}=\beta_k|_{\Gamma_N}=1\}.
 \tag{A.33.1}
\]
At every level, the literal r- and s-frames can therefore
be retained with the actual finite normalizers \(S_3\)
and \(C_2^2\) from A.30. This strengthens the compatible
actual-action assertion without asserting a full
order-two lift.

\begin{proof}
Use the arithmetic \(\sigma\) of A.31 and set
\(u_N=q_N(A_\sigma)\), \(b_N=q_N(\beta_\sigma)\).
We first prove the matrix fact replacing the special
modulo-twelve adjustment in A.30. In signed matrices put
\[
 R=\begin{pmatrix}0&-1\\1&-1\end{pmatrix},\qquad
 D=\begin{pmatrix}0&1\\1&0\end{pmatrix}.
\]
The signed representative \(u_N\) centralizes \(R\)
and is identity modulo four. Indeed projective
centralization cannot replace \(R\) by \(-R\):
their traces are \(-1\) and 1, different modulo
every \(N\) under consideration. The signed
modulo-four assertion is part of R3's pure-frame
normalization. Since \(R\) is cyclic, its matrix
centralizer consists of \(vI+wR\); for a
determinant-one element \(u\) of this centralizer,
\[
              DuD=u^{-1},\qquad
              (u^{-1}D)^2=I,
              \operatorname{tr}(u^{-1}D)=0.
 \tag{A.33.2}
\]

The reflection \(u^{-1}D\) is conjugate to \(D\)
by an element of \(\operatorname{SL}_2(\mathbb Z_p)\)
at every prime, with conjugator identity modulo two
at the dyadic prime. For odd \(p\), its two eigenvalues
are 1 and \(-1\), and its eigensummands are free
rank-one direct summands. Thus a general linear
conjugator exists. The determinant map on
\(C_{\operatorname{GL}_2(\mathbb Z_p)}(D)\) is onto,
so its determinant can be corrected to one.
For \(p=2\), put \(T=u^{-1}D\equiv D\pmod4\)
and take the cyclic-basis matrix
\[
                   H_0=(e_1,Te_1).
\]
It satisfies \(H_0DH_0^{-1}=T\),
\(H_0\equiv I\pmod2\), and
\(d_0=\det H_0\equiv1\pmod4\).
The matrix
\[
 K=\begin{pmatrix}v&w\\w&v\end{pmatrix},\qquad
 v=(d_0^{-1}+1)/2,\quad w=(d_0^{-1}-1)/2,
\]
centralizes \(D\), is identity modulo two,
and has determinant \(d_0^{-1}\).
Consequently \(H_0K\) is the required special
linear conjugator. Reduction and the Chinese
remainder theorem give the finite pure conjugator
at every \(N\). It lifts to a pure element of
\(\widehat B\), because reduction modulo two
is the quotient \(\widehat B/F\cong S_3\).

At the dyadic relative level, retain the full
Galois involution \(a_2=\Phi_\sigma c\).
Its r-normalizing involution
\(D_2=J_2^{-1}a_2\), with
\(J_2=\Phi_\sigma(s)A_{\sigma,2}\), has finite
image \(u_N^{-1}d\), where \(d=sc\).
The matrix calculation provides a pure conjugate
of \(D_2\) with finite image \(d\).
Apply A.16 to the inverse image of
\(\langle q_N(r),d\rangle\cong S_3\), retaining
that involution. Coprime conjugacy of its
order-three complement with \(\langle r\rangle\)
then gives \(D'_2\) inverting \(r\) and having
finite image \(d\). The correction
\(D_2D'_2\) centralizes \(r\) and has finite
image \(u_N^{-1}\). It changes the r-frame
to finite image one while keeping its associated
r-normalizer involutory. At odd primes use the
identity frame and \(sc\), as in A.30.

For the s-frame, \(b_N\) centralizes the signed
matrix \(S\). Conjugation by
\(C=\operatorname{diag}(-1,1)\) sends it to its
inverse, because \(CSC=S^{-1}\) and the
determinant-one \(S\)-centralizer is the
corresponding quadratic norm-one group.
Thus \(b_N^{-1}c\) is an involution in the
common finite extension. At each odd prime the
centralizer epimorphism in A.23, followed by
Schur--Zassenhaus, supplies an involution above
it. At two keep the full Galois s-normalizer.
This gives the same equivariant s-frame gluing
as in A.30.

R2 now supplies a full lift at each \(N\) with
these literal frames. The proof of A.30 shows
that its square has identity actual action and
both relative frames equal to one. Restrict
the nonempty compact sets \(\mathcal C_N\)
from A.31 by these additional closed conditions
on the frames of \((hc)^2\). A common multiple
again provides the finite-intersection property,
so one full \(h\) satisfies the conditions for
all \(N\). This proves (A.33.1).
\end{proof}

The matrix certificate
\path{check_A33_equivariant_tower.py} constructs
the pure special-linear conjugators for every
allowed r-Cartan value at moduli
12, 24, 48, 60, 84, 120, and 132. The proof above
establishes all moduli. In the specified full
framed fibre the remaining equation is
\[
            v(ava^{-1})=a^{-2},
            \qquad v\in L_{\mathrm{cong}}^{\mathrm{fr}}.
\]
This is again a sufficient route to a full
involution, not a claimed solution or a necessary
restriction on differently framed lifts.

\statement{A.34}{Proved}{The full peripheral lifting theorem holds over any closed normal relative kernel}
Let \(\Pi\) be either the full free three-cusp group
\(F\), or the full modular group \(\widehat B\), with
the respective three peripheral elements
\((x,y,z)\) or \((s,r,\tau^{-1})\).
Let \(V\triangleleft\Pi\) be any closed normal subgroup,
without an openness or characteristicity assumption,
and put \(R=R_{\mathrm{nil}}(V)\).
Every continuous automorphism of \(\Pi/R\) preserving
the three positive peripheral conjugacy classes, with
specified literal conjugators there, lifts to a full
peripheral automorphism of \(\Pi\) with those conjugators.
The lift stabilizes \(R\). If a specified peripheral
conjugator is one, that peripheral can be fixed literally.

For a descending directed family of closed subgroups
\((V_i)\) of any profinite group, one also has
\[
 R_{\mathrm{nil}}\Bigl(\bigcap_i V_i\Bigr)
                  =\bigcap_i R_{\mathrm{nil}}(V_i).
 \tag{A.34.1}
\]
In particular, if \(K_{\mathrm{cong}}=\bigcap_NU_N\),
the group kernel \(T\) in A.31 is exactly
\(R_{\mathrm{nil}}(K_{\mathrm{cong}})\).

\begin{proof}
For every profinite group \(V\), with all commutator
subgroups taken closed,
\[
 R_{\mathrm{nil}}(V)=\bigcap_{n\ge1}\gamma_n(V),
 \qquad
 q(R_{\mathrm{nil}}(V))=
                   \gamma_\infty(q(V))
 \tag{A.34.2}
\]
for every finite continuous image \(q(V)\).
The first equality follows because bounded-class
nilpotent profinite quotients are separated by their
finite nilpotent quotients. For the second, a
descending family of nonempty compact fibres of
a fixed finite image value has nonempty intersection.
The finite lower central series stabilizes, so
\[
 q([R,V])=[\gamma_\infty(q(V)),q(V)]
          =\gamma_\infty(q(V))=q(R).
\]
Finite quotients separate closed subgroups. Therefore
\([R,V]=R\), and normality gives
\[
                             [R,\Pi]=R.
 \tag{A.34.3}
\]

Apply the finite scalar-cover and complete proper-fibre
count used in A.3. For a finite quotient
\(q_E:\Pi\twoheadrightarrow E\), take
\(A=q_E(R)\). The auxiliary actual source map
\(q^\dagger:\Pi\twoheadrightarrow D^\dagger\)
has central image of \(R\), because its composite
to \(E/A\) kills \(R\).
Equation (A.34.3) implies that \(q^\dagger\) itself
kills \(R\). Thus the specified automorphism of
\(\Pi/R\) composes with this actual quotient map
to give the required new epimorphism to
\(D^\dagger\), with the prescribed conjugators.
The finite count gives a positive proper fibre over
every joint finite constraint, including all the
specified conjugator coordinates.

Compactness of \(\Pi^3\) now supplies full
conjugators, the peripheral product relation, and
generation in every finite quotient, exactly as
in A.3. The reference D.50 count explicitly applies
to both sources: in the modular case conjugacy
to \(s,r\) supplies their defining orders, so no
homomorphism from the wrong free source is inserted.
The resulting continuous endomorphism of the
finitely generated group \(\Pi\) is surjective
and hence an automorphism. If it is denoted by
\(\Phi\), and \(a\) is the specified quotient
automorphism, then \(q_R\Phi=a q_R\).
Since both maps are automorphisms, this implies
\(\Phi^{-1}(R)=R\) and hence \(\Phi(R)=R\).
Thus preservation is a conclusion, not a missing
characteristicity assumption. If one conjugator
is one in the retained quotient, its full lift
lies in \(R\); the usual inner correction fixes
that peripheral literally without changing the
retained map or any prescribed quotient conjugator.

Finally take any finite image \(q\) of the ambient
group containing the \(V_i\). Compactness gives
\(q(\bigcap_iV_i)=\bigcap_iq(V_i)\).
The descending directed family of subgroups of
this finite group stabilizes. Its lower central
residual consequently commutes with that intersection.
Using (A.34.2) on each subgroup gives equal finite
images for the two sides of (A.34.1). Both sides
are closed, proving equality. Apply this to
the family \(U_N\) to obtain the final assertion.
\end{proof}

The same audit extends A.3's central-quotient criterion
from characteristic to closed normal \(N\), when
``lift'' explicitly means an automorphism inducing
the given automorphism of \(\Pi/N\). Necessity
uses its resulting preservation of \(N\) and
\([N,\Pi]\); sufficiency is the unchanged finite
count. These extensions concern ordinary peripheral
automorphisms. They do not solve an involution
norm equation or construct a desired automorphism
of an unspecified infinite quotient.

\statement{A.35}{Proved}{The full framed congruence correction kernel contains a closed dyadic procyclic subgroup}
Let \(L_{\mathrm{cong}}^{\mathrm{fr}}\leq\mathcal H\)
be the kernel in A.33, retaining the actual action and
both literal full relative frames at every principal
congruence level. Then
\[
             \mathbb Z_2\ \lhook\joinrel\longrightarrow
                    L_{\mathrm{cong}}^{\mathrm{fr}}
 \tag{A.35.1}
\]
is a continuous closed embedding. The image may also
be required to lie in all three corrected resolvent
kernels of B.53. In particular the remaining norm
kernel of A.33 is not pro-odd. This statement alone
neither obstructs nor solves its twisted norm equation.

\begin{proof}
Write
\[
 R=R_{\mathrm{nil}}(K_{\mathrm{cong}})
     =\bigcap_{12\mid N}R_{\mathrm{nil}}(U_N),
       \qquad \Pi=\widehat B.
\]
The equality is A.34. Use the specified tripod map
\(q_0:\widehat K\twoheadrightarrow A_9\) and the
fifteen normal filling kernels \(N_i\) of B.57.
We first check the extra retained kernel:
\[
                         q_0(R)=A_9.
 \tag{A.35.2}
\]
Indeed \(R\) is normal in \(\Pi\) and contained in
\(\widehat K\). Its image is normal in the simple
group \(A_9\). If that image were trivial, \(q_0\)
would factor through \(\widehat K/R\), an open normal
subgroup of \(\Pi/R\). Taking the core of the finite
quotient kernel would make \(A_9\) a composition
factor of a finite quotient of \(\Pi/R\).
Every such quotient factors through one of the relative
congruence groups in A.31. Their nonabelian composition
factors are only \(\operatorname{PSL}_2(\mathbb F_p)\)
for primes \(p\). None is \(A_9\): its order is
\(181440\), and a prime \(p\geq5\) dividing that
order is only \(5\) or \(7\), whose corresponding
simple groups have orders \(60\) and \(168\).
This contradiction proves (A.35.2).

Set
\[
                 V=R\cap\bigcap_{i=1}^{15}N_i.
\]
Each of these sixteen closed normal subgroups of
\(\widehat K\) has full \(q_0\)-image, by (A.35.2)
and B.57. Their iterated closed commutator lies in
\(V\) and still has image \(A_9\), since \(A_9\)
is perfect. Thus \(q_0(V)=A_9\). The intersection
is normal in \(\Pi\). The induced sixty-coordinate
quotient of B.57 consequently again has
\[
 q(V)=M=A_9^{60},\qquad
 \Pi/(V\cap\ker q)=M\rtimes(\Pi/V).
\]
All subsequent central-refinement arguments in B.57
apply with this smaller \(V\). The diagonal odd
transposition on \(M\) extends by the identity on
\(\Pi/V\), and extends to the actual central
refinement while preserving the classes of \(s,r,\tau\).
Moreover the three conjugators constructed there
belong to the perfect subgroup \(M^\#\). They
therefore have image one on \(\Pi/V\).

Apply A.3 and A.34, retaining these projected
conjugators. The full cusp conjugator belongs to
\(V\); conjugating the resulting automorphism by its
inverse normalizes it to fix \(\tau\) literally.
The resulting \(h\in\mathcal H\) still induces the
identity on \(\Pi/V\), and its two unique pure
frames belong to \(V\). Since \(V\subseteq R\),
it induces the identity on \(\Pi/R\) and its two
frames have image one there. In particular
\(h\in L_{\mathrm{cong}}^{\mathrm{fr}}\).
Its three resolvent images are one, because \(V\)
is contained in the original fifteen filling kernels.

The tripod outer action of \(h\) on the same
\(A_9\) quotient is the nontrivial class in
\(\operatorname{Out}(A_9)=C_2\). Cusp normalization
changes that action only by an inner automorphism:
its conjugator lies in \(V\subseteq\widehat K\).
As in B.57, this gives a continuous quotient of order
two of the closed procyclic group
\(\overline{\langle h\rangle}\). R1 says that
\(\mathcal H\) is torsion-free. A torsion-free
procyclic profinite group with a quotient of order
two has dyadic factor \(\mathbb Z_2\), proving
(A.35.1). Closedness of the action and frame
conditions keeps the entire procyclic group in
\(L_{\mathrm{cong}}^{\mathrm{fr}}\).
\end{proof}

The construction supplies no formula for the action
of the chosen negative lift \(a\) on this procyclic
subgroup, nor a value of \(a^2\) in its abelian
quotients. It therefore does not replace the
calculation of \(H^2(C_2,M)\) in A.19. Its exact
consequence for Lead~3 is that a coprime splitting
argument cannot be justified by declaring the
noncongruence correction kernel pro-odd.


\statement{A.36}{Proved}{A peripheral conjugator to a Grothendieck--Teichm\"uller involution satisfies the pentagon}
Let \(\alpha\in\widehat{GT}\) be an involution and let
\(g\in\mathcal P\) satisfy \(\alpha=g c g^{-1}\).
Then \(g\in\widehat{GT}\). Consequently
\[
 \widehat{GT}\cap\{g c g^{-1}:g\in\mathcal P\}
   =\{g c g^{-1}:g\in\widehat{GT}\}.
 \tag{A.36.1}
\]
For an involution already in \(\widehat{GT}\), conjugacy
to \(c\) in any of \(\mathcal P,\widehat{GT}_0,
\widehat{GT}\) is therefore equivalent. This statement
does not assert that every involution is so conjugate.

\begin{proof}
A.6 first puts \(g\) in \(\widehat{GT}_0\). Write
\(\alpha=(-1,f)\) and \(g=(\mu,h)\) in our normalized
convention. The equality \(\alpha g=g c\), computed
on \(y\), and uniqueness of the normalized conjugator
give
\[
                         h=f\,\alpha(h).
 \tag{A.36.2}
\]
We shall use the actual action of \(\alpha\) on the
completed parenthesized braid operad
\(\widehat{\mathrm{PaB}}\). We do not presume that
\(g\) acts on its four-strand component.

Here is the groupoid convention explicitly. Objects are
labelled binary parenthesizations. Composition is
written right to left, and the endomorphism group of
any object with \(n\) labels is \(\widehat P_n\).
Let
\[
 b=((12)3),\qquad b'=(1(23)),\qquad A:b\longrightarrow b'
\]
be the associativity arrow with underlying identity
braid. Let \(x\) be the full twist of strands one and
two at \(b\), and let
\(y=A^{-1}y' A\), where \(y'\) is the full twist of
strands two and three at \(b'\). These give the free
factor \(\widehat F(x,y)\) in
\(\widehat P_3=\widehat{\mathbb Z}\times\widehat F_2\).
Choose the customary source-based convention
\[
             \alpha(A)=A f,
 \qquad \alpha(x)=x^{-1},
 \qquad \alpha(y)=f^{-1}y^{-1}f.
 \tag{A.36.3}
\]
The last identity follows directly by applying
\(\alpha\) to \(A^{-1}y'A\), so this convention
agrees with the normalization at the start of this
part. The associator parameter has no central full-twist
coordinate and belongs to the commutator subgroup;
these are exactly the normalization conditions in the
parenthesized-braid presentation
\cite[Definitions 6.10--6.11, Lemma 7.1 and Proposition 7.3]{CEBHR2019}.

Consider the genuine arrow
\[
                     \theta=A h:b\longrightarrow b'.
\]
It exists for every \(h\in\widehat F(x,y)'\), whether
or not that parameter satisfies the pentagon. Equations
(A.36.2)--(A.36.3) imply
\[
             \alpha(\theta)=A f\alpha(h)=A h=\theta.
 \tag{A.36.4}
\]
Now form its five operadic insertions between the
following actual four-strand objects:
\[
 \begin{aligned}
 b_0&=(((12)3)4),& b_1&=((1(23))4),\\
 b_2&=(1((23)4)),& b_3&=(1(2(34))),\\
 b_4&=((12)(34)).
 \end{aligned}
\]
The three-edge path from \(b_0\) to \(b_3\) consists of
\[
 \begin{aligned}
 e_1&=\theta^{1,2,3}:b_0\to b_1,\qquad
 e_2=\theta^{1,23,4}:b_1\to b_2,\\
 e_3&=\theta^{2,3,4}:b_2\to b_3.
 \end{aligned}
\]
The two-edge path consists of
\[
 d_1=\theta^{12,3,4}:b_0\to b_4,\qquad
 d_2=\theta^{1,2,34}:b_4\to b_3.
\]
The superscripts mean literal cabling by identity arrows
of the indicated binary objects; an unused outer
strand is also inserted by an identity arrow. Since
\(\alpha\) is an object-fixing operad automorphism,
it fixes those identity arrows and commutes with these
insertions. Equation (A.36.4) therefore fixes all five
arrows, not merely their classes. Their discrepancy is
the based pure braid
\[
                D=(d_2d_1)^{-1}(e_3e_2e_1)
                   \in\operatorname{End}(b_0)=\widehat P_4,
 \qquad \alpha(D)=D.
 \tag{A.36.5}
\]

C.28 gives \(\operatorname{Fix}_{\widehat P_4}(\alpha)=1\).
For clarity its application here uses actual strand
forgetting, compatible with the unital operadic action
\cite[Proposition 11.5 and Corollary 11.6]{CEBHR2019}.
The full profinite pure-braid sequences have free
kernels of ranks two and three. Each finite
point-pushing cusp is mapped to a conjugate of its
inverse. The infinity cusp is represented, up to
inverting the peripheral generator, by
\(z_m z_{m-1}^{-1}\), where \(z_j\) is the full
central twist of the first \(j\) strands. Both twists
are sent to their inverse; choose a parenthesization
grouping the first \(m-1\) strands to make this
literal. Thus the kernel action has all the negative
peripheral classes required by C.28. Induction from
\(\widehat P_2=\widehat{\mathbb Z}\), on which the
action is inversion, proves the displayed fixed-subgroup
assertion. In particular (A.36.5) has \(D=1\).

This is precisely the pentagon for the parameter
\(h\). The two hexagons already hold because
\(g\in\widehat{GT}_0\). If the half-braiding arrow is
denoted \(R:(12)\to(21)\), its parameter \(\mu\) is
implemented by the well-typed arrow
\[
 R_\mu=R\,(R_{\mathrm{op}}R)^{(\mu-1)/2}:(12)\to(21),
 \qquad R_{\mathrm{op}}:(21)\to(12),
\]
rather than by an untyped groupoid power. Every unit
\(\mu\in\widehat{\mathbb Z}^{\times}\) is odd.
The parenthesized-braid presentation identifies these
hexagons and this pentagon with the defining
Grothendieck--Teichm\"uller equations
\cite[Lemma 7.1 and Proposition 7.3]{CEBHR2019}.
Since \(g\) already is an invertible normalized
pair in \(\widehat{GT}_0\), the extra pentagon puts
it in \(\widehat{GT}\). More explicitly, the pentagon
solutions inside the profinite group \(\widehat{GT}_0\)
form a closed submonoid by the operadic presentation.
Every closed submonoid of a profinite group is a
subgroup: the powers \(g^{n!-1}\) converge to
\(g^{-1}\). This also checks the unit condition in
the operadic monoid definition of \(\widehat{GT}\).
\end{proof}

The decisive fixed-point argument used only the honest
operadic action of the given \(\alpha\in\widehat{GT}\).
It does not construct a five-point action of an
arbitrary element of \(\widehat{GT}_0\), nor prove
that a general involution there has such an action.
Thus A.36 is a descent of an existing conjugator,
leaving the existence questions in A.13--A.14 open.

\statement{A.37}{Proved}{Every full Grothendieck--Teichm\"uller involution gives a unique pentagon solution}
Let \(\alpha=(-1,f)\in\widehat{GT}\) be an involution.
Let \(k\in F'\) be the unique solution supplied by R1
of \(f=k^{-1}\alpha(k)\), and put \(u=k^{-1}\).
Then \(u\) satisfies the full profinite pentagon.
Moreover \(\alpha\) is conjugate to \(c\) in
\(\widehat{GT}\) if and only if there exists
\(\mu\in\widehat{\mathbb Z}^{\times}\) such that
\[
                 x^\mu u^{-1}y^\mu u\sim_F(xy)^\mu.
 \tag{A.37.1}
\]

\begin{proof}
The equation for \(k\) is equivalent to
\(u=f\alpha(u)\). Therefore the arrow \(A u\)
is fixed by the genuine action of \(\alpha\),
exactly as in (A.36.4). Its five inserted arrows
have a fixed pentagon discrepancy in \(\widehat P_4\),
which C.28 makes one. This proves the pentagon
without assuming that a pair \((\mu,u)\) is an
endomorphism of the operad.

If (A.37.1) holds, the normalized substitution
\(g=(\mu,u)\) is an endomorphism preserving the
three named cusp classes with the common unit
multiplier \(\mu\). C.29 makes it a full
automorphism, so \(g\in\mathcal P\).
The identity \(u=f\alpha(u)\) gives
\(\alpha g=g c\) on both free generators, and
A.36 puts \(g\) in \(\widehat{GT}\).
Conversely any such conjugator \(g=(\mu,v)\)
has \(v=f\alpha(v)\). The uniqueness in R1
implies \(v=u\), and its peripheral property gives
(A.37.1).
\end{proof}

The remaining condition is the full third cusp
with a suitable unit multiplier. Full generation
follows once that equation holds. This automatic
generation and the general peripheral conjugator
criterion already occur in the supplied LM4 D.35;
C.29 independently proves the needed generation
statement by finite-cover boundary homology. The
additional result here is that the canonical \(u\)
already satisfies the pentagon. A.5 shows that the
third cusp need not hold with multiplier one.
The generation condition is independent of \(\mu\):
unit powering preserves each of the two closed
procyclic subgroups. The choice of multiplier
affects the third-cusp condition.
No prounipotent pentagon-to-hexagon implication
has been transferred to the full profinite setting.
Furusho's pro-$\ell$ and pronilpotent conclusions
allow a quadratic extension for the multiplier;
the full profinite implication is posed separately
as Question~14 in that paper
\cite[\S3, Corollaries 12--13 and Question 14]{Furusho2010}.

\statement{A.38}{Proved}{The prescribed coboundary gives peripheral conjugators at two and three}
Let \(\alpha\in\widehat{GT}\) be an involution, and let
\(u\in F'\) be its canonical parameter from A.37.
Use the lower central series
\(\gamma_2(F)=[F,F]\), \(\gamma_3(F)=[\gamma_2(F),F]\)
and define \(t\in\widehat{\mathbb Z}\) by
\[
              u\equiv[x,y]^t\pmod{\gamma_3(F)},
              \qquad D=1+24t.
 \tag{A.38.1}
\]
The commutator convention is the one whose logarithm
has leading term \([X,Y]\), as in A.4.
For a prime \(p\), there exists a peripheral
conjugator from \(c_p\) to \(\alpha_p\) of the
specific normalized form
\[
 g_p(x)=x^{\mu_p},\qquad
 g_p(y)=u_p^{-1}y^{\mu_p}u_p,
 \qquad \mu_p\in\mathbb Z_p^{\times},
 \tag{A.38.2}
\]
if and only if \(D_p\) is a square unit in
\(\mathbb Z_p\). This condition always holds at
\(p=2\) and \(p=3\). For \(p\geq5\), it is equivalent
to \(D_p\bmod p\) being a nonzero quadratic residue.
In particular every involution already in full
\(\widehat{GT}\) is peripherally conjugate to the
standard involution on its maximal pro-\(2\) quotient.

\begin{proof}
A.37 proves the full pentagon for \(u\), hence its
pro-\(p\) pentagon. Furusho's Corollary~12 gives the
two hexagons with parameter
\(\mu_p=\pm\sqrt{1+24t_p}\)
\cite[Theorem 2 and \S3, Corollary 12]{Furusho2010}.
The theorem generally permits a quadratic extension;
we use it only when the chosen root belongs to
\(\mathbb Z_p^{\times}\). Then the half-parameter
\(m_p=(\mu_p-1)/2\) belongs to \(\mathbb Z_p\):
for odd \(p\) this is immediate, and for \(p=2\)
it follows because \(\mu_p\) is odd. Consequently
all words in the hexagons lie in \(F^{(p)}\).
The injective comparison
\(F^{(p)}\hookrightarrow F_2(\mathbb Q_p)\)
used in Furusho's proof therefore makes the
hexagons identities in the actual pro-\(p\) group,
not merely in its prounipotent completion.
They give the three peripheral power classes.
The substitution (A.38.2) has mod-\(p\) homology
matrix \(\mu_p I\), hence is an automorphism by the
pro-\(p\) Frattini criterion. The exact identity
\(u=f\alpha(u)\) from A.37 gives
\(\alpha_p g_p=g_p c_p\) on both free generators.

Conversely, apply the necessary third-cusp coefficient
calculation A.4 in the class-three
\(\mathbb Q_p\)-Malcev quotient. This calculation is
valid for every prime after passing to \(\mathbb Q_p\),
and gives
\(\mu_p^2=1+24t_p\). A unit multiplier therefore
requires a square unit. Finally
\[
       D_2\in1+8\mathbb Z_2,
       \qquad D_3\in1+3\mathbb Z_3.
\]
Every element of \(1+8\mathbb Z_2\) is a square
unit, and Hensel's lemma gives a square root in
\(\mathbb Z_3^{\times}\) of every element of
\(1+3\mathbb Z_3\). For odd \(p\), the unit-square
criterion is exactly the nonzero-square criterion
in the residue field. This proves all assertions.
\end{proof}

The parameter in (A.38.1) uses the lower-central
quotient \(\gamma_2(F)/\gamma_3(F)\), not the
abelianization of the derived group. This is also
the integral coefficient represented by the
\(XY\)-coefficient in Furusho's exponential
notation. The displayed derived-series expression
preceding Question~14 on printed page~555 of that
paper has a typographical error: its asserted
cyclic quotient is the lower-central one used here.

A.38 does not settle the dyadic case for arbitrary
\(\mathcal P\)- or \(\widehat{GT}_0\)-involutions:
the genuine full \(\widehat{GT}\)-action supplied
A.37's pentagon. It also concerns maximal pro-\(p\)
quotients, not all relative quotients \(F/U^{(p)}\).
For other local conjugators, whose parameter need
not be \(u_p\), the necessity assertion about
\(D_p\) is not being made.
For the actual pentagon word supplied by a full
\(\widehat{GT}\)-involution, A.42 below proves
that this square-unit criterion is always satisfied.

\statement{A.39}{Proved}{The audited full conjugacy question is one unit third-cusp equation}
Let \(\alpha=(-1,f)\in\mathcal P\) be a nonidentity
involution. Let \(u\in F'\) be the unique solution
of \(u=f\alpha(u)\), supplied by R1 after inverting
its coboundary parameter. Define
\[
 \mathcal M_\alpha=
 \{\mu\in\widehat{\mathbb Z}^{\times}:
       x^\mu u^{-1}y^\mu u\sim_F(xy)^\mu\}.
 \tag{A.39.1}
\]
Then \(\alpha\) is conjugate to \(c\) in
\(\mathcal P\) if and only if
\(\mathcal M_\alpha\ne\varnothing\).
This set is either empty or exactly
\(\{\mu,-\mu\}\) for one unit \(\mu\).
If \(\alpha\) belongs to \(\widehat{GT}_0\),
respectively \(\widehat{GT}\), both resulting
conjugators belong to that same group.
The existence of a unit satisfying (A.39.1) is
detected by all characteristic finite quotients of
\(F\). For a specified unit \(\mu\), the validity
of (A.39.1) is not detected by the maximal
pro-\(p\) quotients alone, even with \(u=1\).

The basic one-equation criterion and automatic
generation were asserted in the supplied LM4 D.35;
D.49 asserted the two-element conjugator fibres.
This statement audits those conclusions using
C.29's independent covering proof, R1's now-unique
parameter, and A.6. A.36 supplies the additional
descent to full \(\widehat{GT}\).

\begin{proof}
If \(\mu\in\mathcal M_\alpha\), the free universal
property defines the endomorphism
\[
       g_\mu(x)=x^\mu,\qquad
       g_\mu(y)=u^{-1}y^\mu u.
\]
Equation (A.39.1) makes it peripheral at all three
cusps with the common unit power. C.29 makes it
automatically onto, hence an automorphism.
The identity \(u=f\alpha(u)\) gives
\(\alpha g_\mu=g_\mu c\). Thus it is the desired
conjugator in \(\mathcal P\). Conversely every
normalized conjugator has this parameter \(u\)
by its uniqueness in R1 and satisfies (A.39.1).
A.6 and A.36 give the two membership assertions.

If \(\mu,\nu\in\mathcal M_\alpha\), then
\(g_\nu^{-1}g_\mu\in C_{\mathcal P}(c)=\{1,c\}\)
by A.6. Its multiplier gives \(\mu=\nu\) or
\(\mu=-\nu\). Both signs occur because
\(g_\mu c=g_{-\mu}\), and these two units are
distinct. This proves the asserted solution set.

For finite detection, let \(Q=F/V\) range over
characteristic finite quotients, and let
\(e_Q\) be its exponent. Test whether there is
\[
 a\in(\mathbb Z/e_Q\mathbb Z)^{\times}
 \quad\hbox{with}\quad
 \bar x^a\bar u^{-1}\bar y^a\bar u
          \sim_Q(\bar x\bar y)^a.
 \tag{A.39.2}
\]
Every such residue \(a\) lifts to a profinite unit,
and every conjugator in \(Q\) lifts to \(F\).
Thus (A.39.2) is the nonemptiness of a closed subset
of the compact space
\(\widehat{\mathbb Z}^{\times}\times F\), where
we retain the unit and the third-cusp conjugator.
A common characteristic refinement gives the
finite-intersection property. Nonemptiness at
every finite quotient is therefore equivalent
to a full solution, by compactness.

For the final assertion take \(\alpha=c\), so
\(u=1\), and choose the profinite unit with
\[
          \mu_2=-1,\qquad \mu_p=1\quad(p\text{ odd}).
 \tag{A.39.3}
\]
At every odd prime the third-cusp equation is
literal. At two, \(x^{-1}y^{-1}\) is conjugate
to \((xy)^{-1}\), so it holds there as well.
But \(\mu\equiv7\pmod{12}\). In the actual
quotient \(F\twoheadrightarrow S_4\) given by
\[
                x=(234),\qquad y=(1234),
\]
with products acting right to left, direct
multiplication gives
\[
     x^7y^7=(12),\qquad (xy)^7=(1423).
 \tag{A.39.4}
\]
A transposition is not conjugate to a four-cycle.
The script \path{check_A39_mixedpower.py} verifies
that the displayed pair generates all twenty-four
elements and checks the full conjugacy class.
This is a finite mixed-prime obstruction to the
particular multiplier (A.39.3), despite every
maximal pro-\(p\) equation holding.
\end{proof}

For \(\alpha\in\widehat{GT}\), A.37 adds the
pentagon for its prescribed \(u\). The later
A.42 proves that full unit roots of
\(\mu^2=1+24c_2(u)\) exist and that every local
hexagon holds. The remaining question is whether
\emph{some} full root gives the full hexagon.
This existential assertion would settle conjugacy:
the two-cycle already follows from the pentagon,
the three-cycle gives the third cusp by C.30,
and C.29 supplies full generation. C.34 explains
why the literal prescribed-root version of
Furusho's Question~14 fails: the mixed-sign
example (A.39.3)--(A.39.4) also fails its full
hexagon. It does not refute existence of a
suitable root, since the same example has the
two standard roots \(1,-1\).

\statement{A.40}{Proved}{A parity-one involution survives an actual noncongruence perfect-socle refinement}
Put
\[
 \Pi=\widehat B,\qquad
 R=R_{\mathrm{nil}}(K_{\mathrm{cong}}),\qquad C=\Pi/R,
 \qquad \widehat K=\ker(\Pi\longrightarrow A_5).
\]
There is a finite quotient
\[
 q:\Pi\twoheadrightarrow Q=M\rtimes A_5,
 \qquad M=A_9^t\quad(t\geq60),
 \tag{A.40.1}
\]
whose kernel is invariant under the full
\(\widehat{GT}_0\)-action, and for which the
\(A_5\)-action on the simple factors is a union of
regular orbits. Let
\[
 W=R\cap\ker q,
 \qquad D=\Pi/W,
 \qquad F_D=F/W.
\]
There is a full \(a\in\widehat{GT}_0\) with
\(\lambda(a)=-1\) and \(\Delta(a)=1\) whose actual
induced automorphism on \(D\) has order two.
Its outer image on \(F_D\) has order two and
projects to the entire compatible family of A.31.
The literal \(C\)-frames from A.33 may also be retained.

This quotient genuinely extends the relative
principal-congruence tower: \(D\) maps onto \(C\)
and has an \(A_9\) composition factor, which no
finite quotient of \(C\) has. The statement does
not assert \(a^2=1\) in the full group, nor that
\(W\) is characteristic under all automorphisms
of the underlying free group.

\begin{proof}
\emph{The finite invariant family.}\par
The twelve cusps of \(\widehat K\) are the
\(A_5/C_5\) cosets. The antipodal relation is
intrinsically given by the nontrivial element of
\(N_{A_5}(C_5)/C_5\). Consider all kernels of
epimorphisms \(\widehat K\twoheadrightarrow A_9\)
which factor through a three-cusp filling, whose
retained triple consists of an antipodal pair and
one further cusp, and which send each retained
positive meridian to a seven-cycle. This is a finite
family because \(\widehat K\) is finitely generated.
It is nonempty by the exact generating triple
(B.57.5)--(B.57.6).

Congruence covariance, as used in A.31, preserves
\(\widehat K\). A peripheral automorphism permutes
its cusps and preserves the stated normalizer
description of antipodes. Its common unit powers
also preserve seven-cycle orders. Precomposition
therefore permutes the entire defined family of
kernels. Each kernel has trivial \(A_5\)-stabilizer:
its stabilizer preserves its nonzero cusp support,
and such a triple has trivial stabilizer by B.57.
Thus every orbit has sixty elements.

Choose an epimorphism for one kernel in each orbit,
and induce it from \(\widehat K\) to \(\Pi\) by
the regular wreath-product construction of B.57.
Take the product of these induced maps over their
common \(A_5\)-quotient. The coordinate kernels
on \(\widehat K\) are exactly the distinct kernels
in the family. Their image is the full product:
for a chosen coordinate, every other distinct
maximal normal kernel maps onto its simple target;
an iterated commutator of those other kernels gives
that factor alone. This proves (A.40.1), including
the asserted regular action and invariance of its
kernel.

Every coordinate sends \(R\) onto \(A_9\).
Otherwise it would factor through
\(\widehat K/R\), an open normal subgroup of
\(C\), making \(A_9\) a composition factor of a
finite quotient of \(C\). This contradicts
A.31 and A.35. Normality of \(q(R)\) in the
product \(M\), together with surjectivity onto
each simple factor, gives
\[
 q(R)=M,\qquad
 D=Q\times_{A_5}C\cong M\rtimes C.
 \tag{A.40.2}
\]
For the second equality, lift a prescribed
\(C\)-coordinate to \(\Pi\), then correct its
finite \(M\)-coordinate by an element of \(R\).
The action of \(C\) on \(M\) factors through
\(A_5\).

\emph{The actual central refinement.}\par
Set
\[
 E^\#=\Pi/[W,\Pi],\qquad Z=W/[W,\Pi],
 \qquad R^\#=R/[W,\Pi].
\]
Thus \(Z\) is central in \(E^\#\), and
\(R^\#\) is the full preimage of \(M\).
The canonical image \(M^\#\) of the universal
central cover \((2.A_9)^t\) in \(R^\#\) is
finite and perfect, maps onto \(M\), and is
normal in \(E^\#\). The universal-cover input is
the same primary Schur-multiplier datum used in
B.57.\cite{CEAtlasA9}
Uniqueness of lifts from a perfect group makes
this canonical image invariant under conjugation.
Centrality gives \(R^\#=M^\# Z\). A.34 gives
\(R=[R,\Pi]\), and consequently
\[
 R^\#=[R^\#,E^\#]=[M^\# Z,E^\#]\subseteq M^\#.
\]
Hence
\[
                   R^\#=M^\#,
                   \qquad Z\subseteq M^\#.
 \tag{A.40.3}
\]

Let \(C^\#\) be the preimage in \(E^\#\) of
the specified section \(C\leq M\rtimes C\).
In \(D\), the actual elements \(s,r\) are
\(M\)-conjugate to the pure elements
\(s_C,r_C\) of this section. Indeed, in each
regular wreath orbit their permutations consist
only of two-cycles, respectively three-cycles.
Every base cycle product is one, since the actual
elements have orders two and three. Solving the
coordinate conjugator equations therefore removes
their entire base components, as in B.57.

Lift these conjugators to \(M^\#\). The
corresponding conjugates \(S,R_3\) of the actual
\(s,r\) in \(E^\#\) lie in \(C^\#\), have
orders dividing two and three, and project to
the named \(s_C,r_C\). The universal property
of \(\Pi=C_2\amalg C_3\) gives a homomorphism
\(j:\Pi\to C^\#\) sending \(s,r\) to these
elements. Its projection to \(C\) is the
canonical quotient map. Thus \(j(R)\subseteq Z\)
is central. Since \(R=[R,\Pi]\), it follows
that \(j(R)=1\). The map factors to a continuous
section \(i:C\to C^\#\), proving
\[
                         E^\#=M^\#\rtimes C.
 \tag{A.40.4}
\]

The action of \(C\) on \(M^\#\) factors
through \(A_5\). On \(M\), it is the specified
factor permutation. An automorphism of the
perfect group \(M^\#\) inducing the identity
on its central quotient \(M\) is the identity:
its difference is a homomorphism from \(M^\#\)
to the central abelian kernel. This proves the
assertion for every element in the kernel of
\(C\to A_5\), and also proves uniqueness of the
lifted permutation action. Put
\[
 Q^\#=M^\#\rtimes A_5,
 \qquad E^\#=Q^\#\times_{A_5}C.
 \tag{A.40.5}
\]
For this chosen section the images of the actual
\(s,r\) in \(Q^\#\) are \(M^\#\)-conjugate
to their pure permutations: the section was built
from precisely those chosen conjugates.

\emph{A positive lift retaining the whole congruence quotient.}\par
Let \(\beta_C\) be A.31's actual order-two
automorphism of \(C\), with the literal frames
of A.33. It literally inverts \(\tau_C\), and
its induced action on \(A_5\) is the standard
one. Write \(c_C\) for genuine standard
conjugation and put
\[
                          h_C=\beta_Cc_C.
\]
The map \(h_C\) acts identically on \(A_5\),
fixes \(\tau_C\), and preserves the positive
\(s,r\) classes with the prescribed pure frames.
The formula
\[
 \psi^\#=\operatorname{id}_{M^\#}\rtimes h_C
        =(\operatorname{id}_{Q^\#},h_C)
       \quad\text{on }E^\#=Q^\#\times_{A_5}C
 \tag{A.40.6}
\]
is a continuous automorphism. It fixes the actual
\(\tau\) in \(E^\#\) literally: its finite
\(Q^\#\)-coordinate is unchanged and its
\(C\)-coordinate \(\tau_C\) is fixed. No
assumption about the order of a finite cusp lift
is needed.

To check \(s\), write \(A=q^\#(s)\) and let
\(d_s\in C\) be its specified positive frame,
so that \(h_C(s_C)=d_s s_Cd_s^{-1}\).
Its \(A_5\)-image centralizes \(\bar s\),
because \(h_C\) acts identically on \(A_5\).
The map
\[
          C_{Q^\#}(A)\longrightarrow C_{A_5}(\bar s)
 \tag{A.40.7}
\]
is onto. Conjugate \(A\) to its pure permutation
by the \(M^\#\)-conjugator above; every pure
permutation in \(C_{A_5}(\bar s)\) then
centralizes it, and conjugating back preserves
the \(A_5\)-coordinate. Choose a lift
\(e_s\in C_{Q^\#}(A)\) of the image of
\(d_s\). The pair \((e_s,d_s)\) lies in the
actual group \(E^\#\) and conjugates \(s\)
to \(\psi^\#(s)\), retaining exactly the
specified \(C\)-frame. The identical argument
applies to \(r\).

We have therefore constructed an automorphism of
the actual central refinement
\(\Pi/[W,\Pi]\) preserving all three positive
peripheral classes with literal conjugators.
Apply the closed-normal criterion A.3/A.34.
It gives a full positive peripheral lift inducing
\[
             \psi_D=(\operatorname{id}_Q,h_C)
                \quad\text{on }D=Q\times_{A_5}C,
 \tag{A.40.8}
\]
and retaining the specified quotient conjugators.
Its full cusp conjugator lies in \(W\). An inner
correction by its inverse makes \(\tau\) literally
fixed without changing (A.40.8) or the \(C\)-frames.
Because \(W\subseteq R\subseteq U_{12}\), the
retained pure-frame conditions also ensure purity
of these full frames. The standard modular
identification gives a full \(h\in\mathcal H\).

Finally set \(a=hc\). The kernel of \(q\) and
\(R\) are invariant under genuine \(c\), whose
action on \(Q\) is denoted by \(c_Q\). Thus
\[
       a_D=(c_Q,h_Cc_C)=(c_Q,\beta_C),
       \qquad a_D^2=1.
 \tag{A.40.9}
\]
Both coordinates induce the same action on
\(A_5\), so this is an actual fiber-product
automorphism. It has order exactly two, for
example from its nontrivial inversion on the
order-three quotient of \(\Pi^{\mathrm{ab}}\).
Its outer action on \(F_D\) also has order two,
since its multiplier is \(-1\). Its projection
at level twelve is the one in A.31, so A.24
gives \(\Delta(a)=1\). The character descends
to this actual outer image by projection to
\(G_{12}\); hence this involution is not
conjugate to the standard one in that image.
The full lift's square has only been shown to
vanish on \(D\), as asserted.
\end{proof}

\statement{A.41}{Proved}{The parity-one involution is compatible through a whole semisimple tripod tower}
Keep \(R,C,\widehat K\) from A.40. Let
\(\mathscr S\) range over finite nonabelian simple
groups not isomorphic to
\(\operatorname{PSL}_2(\mathbb F_p)\) for any
prime \(p\). For each such group take all kernels
of epimorphisms from the antipodal-pair tripod
fillings of \(\widehat K\) onto that group.
For every finite collection of targets let
\(q_J:\Pi\twoheadrightarrow Q_J\) be the joint
induced quotient over \(A_5\), and put
\[
 W_J=R\cap\ker q_J,
 \qquad W_* =\bigcap_J W_J,
 \qquad D_* =\Pi/W_*.
 \tag{A.41.1}
\]
There is one full \(a\in\widehat{GT}_0\) with
\(\lambda(a)=-1\) and \(\Delta(a)=1\) whose
actual action on \(D_*\) has order two.
It retains A.33's literal \(C\)-frames and
induces standard conjugation on every \(Q_J\).
This is an assertion about the specified joint
tower; no cofinality or full order-two lift is
asserted.

\begin{proof}
For each fixed finite target there are finitely
many kernels, and the whole family is invariant
under \(\widehat{GT}_0\). A retained cusp cannot
map trivially: killing one of the three cusps
would make the remaining two-cusp quotient
procyclic, which cannot surject onto a nonabelian
simple group. Thus the support is exactly the
retained triple, so every \(A_5\)-orbit is free.
The full-product argument in A.40 now gives
\[
 Q_J=M_J\rtimes A_5,
 \qquad M_J=\prod_{i\in I_J}S_i,
 \qquad q_J(R)=M_J.
 \tag{A.41.2}
\]
The last equality follows as before because none
of the \(S_i\) occurs as a composition factor
of a finite quotient of \(C\).

The only additional input for repeating A.40's
central-refinement argument is finiteness of the
universal central cover of \(M_J\). Here is the
elementary finiteness check. The integral bar
complex of a finite group has finite rank in
each degree, so its second homology is finitely
generated. Transfer from the trivial subgroup
shows that multiplication by the group order
annihilates this homology.\cite[Section~2 and
Corollary~3.4]{ABrown2010} It is therefore finite.
For a perfect group the kernel of its universal
central cover is precisely this second homology,
by the free-presentation construction and Hopf's
formula. The cover is consequently finite and
perfect. Thus every step of (A.40.3)--(A.40.8)
applies to (A.41.2), with no restriction on the
prime divisors of its Schur multiplier.

For a finite collection \(J\), let
\(\mathscr C_J\subseteq\mathcal H\) consist of
the full normalized positive elements whose
actual action on \(\Pi/W_J\) is
\((\operatorname{id}_{Q_J},h_C)\) and whose
two literal \(C\)-frames are the prescribed
ones. These sets are nonempty by the preceding
lifting argument. They are closed: quotient
actions and the normalized frame coordinates
are continuous, and equality of an infinite
quotient action or frame is the intersection
of its closed finite-coordinate equalities,
as in A.31--A.33. For finitely many \(J\),
their union is a single finite collection whose
solution set lies in the intersection. Compactness
of \(\mathcal H\) therefore gives
\[
                         h\in\bigcap_J\mathscr C_J.
\]
Put \(a=hc\). Its action on each joint quotient
is \((c_{Q_J},\beta_C)\), of order two.
These are the actual images of one element,
so their compatibility does not depend on a
surjectivity assertion for torsion transition
maps. Passing to the inverse limit proves the
claim. The image remains nontrivial and of
parity one already at level twelve.
\end{proof}

The finite-family argument controls direct products
of simple factors in free \(A_5\)-orbits, including
their actual scalar central refinements. It does
not control arbitrary nonsimple chief layers of
the remaining kernel. In particular it neither
annihilates every \(H^2(C_2,M)\) in A.19 nor
solves the full nonabelian norm equation.

\statement{A.42}{Proved}{Every full Grothendieck--Teichm\"uller involution has its prescribed conjugator at every prime}
Let \(\alpha=(-1,f)\in\widehat{GT}\) be a
nonidentity involution, and let \(u\in F'\) be
the unique full solution of \(u=f\alpha(u)\)
from A.37. Then
\[
 D=1+24c_2(u)\in
          (\widehat{\mathbb Z}^{\times})^2.
 \tag{A.42.1}
\]
In particular there is a full unit
\(\mu\in\widehat{\mathbb Z}^{\times}\) with
\(\mu^2=D\). For every prime \(p\), and either
local root \(\mu_p\), the prescribed pair
\((\mu_p,u_p)\) belongs to the pro-\(p\)
Grothendieck--Teichm\"uller group and its
normalized peripheral automorphism conjugates
\(c_p\) to \(\alpha_p\).

Any choice of a full square root in (A.42.1)
therefore gives a compatible conjugator on the
maximal pronilpotent quotient
\(\prod_p F_2^{(p)}\). No corresponding full
profinite conjugator or arbitrary relative-level
conjugator is asserted. In particular the statement
does not settle the dyadic case for an arbitrary
\(\mathcal P\)- or \(\widehat{GT}_0\)-involution.

\begin{proof}
A.37 proves the full pentagon for \(u\), so its
image \(u_p\) satisfies the pentagon in the
actual pro-\(p\) pure braid group. C.33 gives
\[
              1+24c_2(u_p)\in(\mathbb Z_p^\times)^2
                  \quad\text{for every }p.
\]
The coefficient is compatible with these quotient
maps. Since
\(\widehat{\mathbb Z}=\prod_p\mathbb Z_p\),
choosing a local square root at each prime yields
(A.42.1) and the stated full unit.

C.33 also gives the actual pro-\(p\) hexagons
for either root, and pro-\(p\) Frattini gives
an automorphism. Equivalently A.38 now applies
at every prime. The full identity
\(u=f\alpha(u)\) projects to the same
intertwining identity, so
\[
       g_p(x)=x^{\mu_p},\qquad
       g_p(y)=u_p^{-1}y^{\mu_p}u_p,
       \qquad \alpha_p g_p=g_p c_p.
\]
Their product is a continuous automorphism of
the maximal pronilpotent quotient and has the
claimed peripheral action and intertwining
property. This proves exactly the asserted
quotient conclusions.
\end{proof}

The additional input over A.38 is C.33's integral
metabelian hexagon identity, whose first possible
\(p\)-denominator excludes nonzero nonsquare
and nonunit parameters, together with C.32's
relative braid-torsion argument excluding zero.
Thus no local square-root obstruction remains
for this canonical word. The remaining full
condition is the third-cusp equation of A.39,
or equivalently the full hexagon for an appropriate
root. A positive answer to Furusho's Question~14
in its \emph{existential-root} variant for this
prescribed \(u\) would now suffice without a
separate square-root assumption. The literal
prescribed-root implication is refuted in C.34;
these two quantifiers must be distinguished.
The independently chosen local signs still give
no full identity: the exact \(S_4\) example
in A.39 already demonstrates that distinction.

\statement{A.43}{Proved}{The actual first scalar central refinement also carries the compatible parity-one involution}
For any finite family in A.40 or A.41, put
\[
       W_1=[W,\Pi],\qquad W_2=[W_1,\Pi].
\]
Then
\[
                         W_2=W_1.
 \tag{A.43.1}
\]
There is a full negative element of parity one
whose actual action on \(\Pi/W_1\) has order
two and retains the prescribed \(C\)-frames.
These refined involutions may be chosen as the
images of one full element for all finite
families simultaneously. Thus A.41 extends to
the actual quotient
\[
            \Pi\Big/\bigcap_J[W_J,\Pi].
 \tag{A.43.2}
\]
The square is still not shown to vanish in the
full group.

\begin{proof}
Write
\[
 E_1=\Pi/W_1=M_1\rtimes C,
 \qquad M_1=R/W_1,\qquad Z_1=W/W_1.
\]
A.40 shows that \(M_1\) is finite and perfect,
\(Z_1\) is central in \(E_1\), and the
\(C\)-action on \(M_1\) factors through
\(A_5\). Since \(M_1/Z_1=M\) is centerless,
\(Z(M_1)=Z_1\).

Consider \(E_2=\Pi/W_2\) and \(M_2=R/W_2\).
The map \(M_2\to M_1\) is a central extension
with kernel \(A=W_1/W_2\), central in all of
\(E_2\). The canonical finite perfect image
\(P\) of the universal central cover of
\(M_1\) satisfies \(M_2=PA\) and is normal
in \(E_2\). Since \(R=[R,\Pi]\),
\[
                  M_2=[M_2,E_2]\subseteq P.
\]
Thus \(M_2=P\) is finite and perfect.

The preimage \(Z_0=W/W_2\) of \(Z_1\) is
central in \(M_2\). Indeed for an element of
that preimage its commutator map on \(M_2\)
takes values in the central kernel \(A\)
and is a homomorphism; perfectness makes it
trivial. Hence \(Z_0\) is finite abelian, with
\[
                    0\longrightarrow A\longrightarrow Z_0
                         \longrightarrow Z_1\longrightarrow0.
 \tag{A.43.3}
\]

Conjugation on \(Z_0\) factors through \(C\),
and in fact through \(A_5\). To check the
second assertion, let \(v\in\ker(C\to A_5)\)
and lift its pure section element in \(E_1\)
to \(E_2\). Its conjugation on \(M_2\)
induces the identity on \(M_1\). An
automorphism of the perfect central extension
\(M_2\) inducing the identity on \(M_1\)
is the identity, by the same central-homomorphism
argument. This proves the asserted factorization.

The resulting \(A_5\)-action on (A.43.3)
is the identity on \(A\) and on \(Z_1\),
because these are central in \(E_2\) and
\(E_1\), respectively. The automorphisms of
this abelian extension which act identically
on both subquotients form the abelian group
\(\operatorname{Hom}(Z_1,A)\): their difference
from the identity is precisely such a homomorphism.
Perfectness of \(A_5\) therefore makes its
action trivial. Thus \(Z_0\) is central in
all of \(E_2\), so the image of
\([W,\Pi]=W_1\) in \(E_2\) is trivial.
This proves \(W_1\subseteq W_2\); the reverse
inclusion is automatic, giving (A.43.1).

Next use A.40's positive automorphism
\(\psi_1=\operatorname{id}_{M_1}\rtimes h_C\)
of \(E_1\), and let \(c_1\) be the genuine
standard action. Put \(b_1=\psi_1c_1\).
Its square is the identity on \(M_1\), since
\(\psi_1\) fixes that subgroup pointwise, and
on \(C\), since \(h_Cc_C=\beta_C\) is an
involution. Both \(b_1\) and \(c_1\) invert
the actual cusp, so \(b_1^2\) fixes \(\tau\).

An automorphism \(\theta\) of \(E_1\) which
fixes \(M_1\) pointwise and induces the identity
on \(C\) differs from the identity by a
homomorphism \(d:C\to Z_1\). Indeed, compare
\(\theta(i(v))\) and \(i(v)\) for the chosen
section. Their ratio centralizes \(M_1\),
hence lies in \(Z_1\), which is central in
all of \(E_1\). The multiplication rule then
makes the ratio an ordinary homomorphism.
But
\[
                  C^{\mathrm{ab}}=\Pi^{\mathrm{ab}}=C_6
\]
is generated by the image of \(\tau\), because
\(R=[R,\Pi]\). If \(\theta\) also fixes the
actual cusp, this homomorphism is trivial.
Apply the observation to \(\theta=b_1^2\)
to obtain \(b_1^2=1\) in the actual refined
quotient.

By (A.43.1), the central refinement required by
A.3 for \(\Pi/W_1\) is \(\Pi/W_1\) itself.
Thus its positive peripheral automorphism
\(\psi_1\) passes the exact lifting criterion
directly. Normalize the full ordinary lift to
fix the cusp and retain its \(C\)-frames, then
multiply by genuine \(c\). This gives the
asserted full element with refined involutory
image and parity one.

Finally \(\psi_1\) is uniquely characterized
by fixing \(M_1\) pointwise, inducing \(h_C\)
on \(C\), and fixing the actual cusp: the
ratio of two such maps is the same central
homomorphism just considered. On enlarging a
finite family the higher refined action
therefore projects to this unique lower one.
Alternatively the nonempty closed full-lift
sets have the finite-intersection property
under finite unions, exactly as in A.41.
Compactness produces one full element with
all the refined images, proving (A.43.2).
\end{proof}

A.43 eliminates every further scalar central
refinement of these particular finite perfect
families. Its argument uses the actual center
and the perfect \(A_5\)-action. A general
noncentral elementary-abelian layer need not
have this structure; its \(C_2\)-norm class
has not been computed here.

\statement{A.44}{Proved}{The canonical substitution passes all dihedral tests and the parity-relative pro-prime generation test}
Let \(\alpha=(-1,f)\in\mathcal P\) be a nonidentity
involution and let \(u\in F'\) be its canonical
parameter, \(u=f\alpha(u)\). Put
\[
 y'=u^{-1}yu,\qquad \psi(x)=x,\quad\psi(y)=y',
 \qquad V=\ker\bigl(F\to C_2,\ x,y\mapsto-1\bigr).
\]
The substitution \(\psi\) is onto every finite
dihedral quotient in which \(x,y\) are generating
reflections. For every prime \(p\), it induces
an automorphism of
\[
                         G_p=F/O^p(V),
 \tag{A.44.1}
\]
and that automorphism conjugates the induced
standard action to the induced \(\alpha\)-action.
This is a generation and unmarked automorphism
statement. It does not assert preservation of
the third cusp by \(\psi\).

\begin{proof}
Write \(E=F\rtimes\langle c\rangle\), where
the named \(c\) acts by \(\alpha\). The three
full section classes have representatives
\[
                         C=c,\qquad X=xc,\qquad Y=y'c.
\]
They are distinct by their mod-two abelian labels
\(0,x,y\). The first two cusp edges have endpoints
\(C,X\) and \(C,Y\). C.28's three-cusp cycle
therefore makes the third cusp's endpoints
\(X,Y\). Every section centralizer in \(F\)
is trivial.

\emph{The induced dihedral action.}\par
For an odd prime \(p\), use
\[
 \begin{gathered}
 q:F\twoheadrightarrow D_{2p},\qquad
 q(x)=a,\quad q(y)=at,\\
 a^2=t^p=1,\quad ata=t^{-1},\qquad
 q(f)=t^\ell,\quad q(u)=t^k.
 \end{gathered}
\]
Then \(q\alpha\) sends \(x,y\) to
\(a,at^{2\ell+1}\). This pair generates because
\(\alpha\) is an automorphism. Every generating
reflection pair of this form is obtained from the
named pair by an automorphism of \(D_{2p}\)
fixing \(a\). Thus \(\ker q\) is
\(\alpha\)-invariant, and the induced map sends
\(t\) to \(t^{2\ell+1}\). Its square is one,
so \(2\ell+1=1\) or \(-1\) modulo \(p\).
The substituted pair fails to generate exactly when
\[
                   q(y')=at^{2k+1}=a.
 \tag{A.44.2}
\]

If the induced action inverts \(t\), then
\(\ell=-1\), and the coboundary equation gives
\(2k=\ell=-1\). The induced action is
\(\operatorname{Int}(a)\). Under the isomorphism
\[
 D_{2p}\rtimes_{\operatorname{Int}(a)}C_2
       \longrightarrow D_{2p}\times C_2,\qquad
 gc^\epsilon\longmapsto(ga^\epsilon,\epsilon),
\]
both \(X\) and \(Y\) map to the central
involution \((1,1)\). Every conjugate of either
class has the same image. The third cusp's two
reflections would therefore have product one in
this quotient, although their product is a
conjugate of \(z\) or \(z^{-1}\), whose image
has order \(p\). This is impossible. Hence
\(\alpha\) acts identically on \(D_{2p}\)
and \(q(f)=1\).

\emph{The pure negative degree-\(p\) cover.}\par
Suppose (A.44.2) still holds and put
\(U=q^{-1}(\langle a\rangle)\). It is
\(\alpha\)-stable. Label \(U\backslash F\)
by \(j\in\mathbb F_p\), using \(Ut^j\).
The right actions of \(x,y,z\) are respectively
\[
                     j\mapsto-j,\qquad j\mapsto1-j,
                             \qquad j\mapsto j-1.
 \tag{A.44.3}
\]
Above each of \(x,y\) there is one width-one
cusp and \((p-1)/2\) width-two cusps. Above
\(z\) there is one width-\(p\) cusp. The
cover has \(p+2\) cusps and Euler characteristic
\(-p\), hence genus zero. Thus \(U\) is the
full profinite sphere group of rank \(p+1\).

The automorphism preserves every \(x\)- and
\(y\)-cusp individually, since its finite action
is identity and their finite frames are
\(q(1)=q(f)=1\). The \(z\)-cusp is unique,
so is preserved regardless of its frame. Every
meridian is inverted up to \(U\)-conjugacy.
Consequently C.28 applies to this pure negative
action on the \((p+2)\)-cusp group.

For clarity, when \(p=3\) the five cusp
representatives can be taken as
\[
        x,\quad yx^2y^{-1},\quad u^{-1}yu,\quad y^2,\quad z^3.
 \tag{A.44.4}
\]
The third is \(U\)-conjugate to the discrete
representative \(yxyx^{-1}y^{-1}\). Their
respective real reflections are \(c\),
\(ycy^{-1}\), \(c\), \(ucu^{-1}\), and
a local \(z\)-reflection adjusted by an even
power of \(z\) to have finite component \(a\).
They all lie in \(U\rtimes\langle c\rangle\).
The width-two frames are
\(y\alpha(y)^{-1}\) and \(f\), both with
trivial finite image; the two width-one frames
are one. The final reflection is \(v_zc\)
with \(q(v_z)=a\), giving an actual \(U\)-frame
for the unique third-cusp lift. No triviality
of that frame is assumed.

\emph{The impossible incidence graph.}\par
The full \(C\)-class splits into precisely
\(p\) \(U\)-classes, denoted \(C_j\).
They are represented by \(gcg^{-1}\) for
\(Ug=Ut^j\). Every such conjugate lies in
the subgroup, and two have the same
\(U\)-class exactly when their left cosets
agree, because \(C_F(c)=1\).
Each of the full \(X,Y\)-classes contributes
exactly one \(U\)-class: its finite component
is \(a\), and a conjugate belongs to the
subgroup precisely when its conjugator has
image in \(C_{D_{2p}}(a)=\langle a\rangle\).
The section vertices are therefore exactly
\(C_j\) together with \(X,Y\).

The width-one \(x\)-edge joins \(C_0\) to
\(X\); the width-one \(y\)-edge, represented
by \(y'\), joins \(C_0\) to \(Y\).
The unique \(z\)-edge has odd width \(p\),
so its endpoints remain in the opposite full
classes \(X,Y\), each of which has just one
\(U\)-class. An even power of \(z\) adjusts
a base reflection into the subgroup, and its
product with \(z^p\) gives the other endpoint.
These edges form a triangle.

For every nonzero pair \(\{j,-j\}\), the
width-two \(x\)-edge joins \(C_j,C_{-j}\).
Its exact local identity is
\(x^2c=xcx^{-1}\): the two conjugators are
\(g\) and \(gx\). For \(y\), use its
actual real reflection \(ucu^{-1}\), which
inverts \(y\). The two endpoints above a
width-two orbit have conjugators \(gu\)
and \(gyu\). Their labels are
\[
                         j+k,\qquad 1-j+k.
\]
By (A.44.2), these are negatives of each other.
The fixed \(y\)-orbit becomes label zero;
every other orbit gives exactly the same
nonzero pair as an \(x\)-edge. Thus the
incidence graph is one triangle together with
\((p-1)/2\) disjoint digons. C.28 requires
one cycle through every edge and vertex, a
contradiction.

This excludes (A.44.2). Failure of generation
in any finite dihedral reflection quotient
would mean that the odd number \(2k+1\)
is not a unit modulo its rotation order.
That failure is detected at an odd prime,
which has just been excluded. This proves
the first assertion.

\emph{The relative Frattini consequence.}\par
A Schreier basis of \(V\) is
\[
                         A=x^2,\qquad B=yx^{-1},\qquad C=xy.
\]
On its abelianization conjugation by \(x\)
sends \(A\) to \(A\), \(B\) to \(C-A\),
and \(C\) to \(A+B\). For odd \(p\),
the invariant plane is spanned by \(A,B+C\),
and the sign line by \(A+B-C\). These three
vectors form a basis, with determinant \(-2\).
Here \(B+C\) is the class of \(y^2\).

The substitution preserves \(V\) and
\(O^p(V)\), so induces an endomorphism of
\(G_p\). Its closed image \(H\) is generated
by \(x,y'\). The intersection with the free
pro-\(p\) kernel \(V^{(p)}\) contains
\(x^2,(y')^2\), whose Frattini images are
the two invariant basis vectors: \(u\in V\)
does not change the abelianized square.
Quotienting by that invariant plane gives
\(D_{2p}\), with \(x,y\) reflections and
\(xy\) having nonzero sign coordinate.
The dihedral result makes \(H\) onto this
quotient, so \(H\cap V^{(p)}\) also covers
the sign line. It therefore covers the entire
Frattini quotient. Frattini gives
\(H\cap V^{(p)}=V^{(p)}\), and hence
\(H=G_p\). Surjectivity makes the induced
endomorphism an automorphism.

At \(p=2\), (A.44.1) is the maximal pro-two
quotient of \(F\), since its kernel and
quotient are pro-two. Its abelianized
substitution is the identity, so ordinary
pro-two Frattini gives the same conclusion.
Finally the full identity \(u=f\alpha(u)\)
gives \(\alpha\psi=\psi c\) on the two
generators and therefore on every stated
quotient. This proves the conjugacy assertion
with precisely its stated unmarked scope.
\end{proof}

\statement{A.45}{Proved}{The canonical word determines the only possible full multiplier}
For \(\alpha,u\) as in A.44, use the universal
full dihedral quotient
\[
 q:F\twoheadrightarrow D_\infty^{\wedge}
      =\widehat{\mathbb Z}\rtimes C_2,\qquad
 q(x)=a,\quad q(y)=at,\quad ata=t^{-1}.
\]
Write \(q(u)=t^k\) and define
\[
                          d_\alpha=2k+1.
 \tag{A.45.1}
\]
Then \(k\in2\widehat{\mathbb Z}\),
\(d_\alpha\in\widehat{\mathbb Z}^{\times}\),
and \((d_\alpha)_2\equiv1\pmod4\).
The involution acts identically on
\(D_\infty^{\wedge}\). Its full conjugacy
criterion is the single equation
\[
 x^{d_\alpha}u^{-1}y^{d_\alpha}u
                   \sim_F(xy)^{d_\alpha}.
 \tag{A.45.2}
\]
If a full peripheral conjugator exists, its
multiplier is exactly \(d_\alpha\) or
\(-d_\alpha\). The same equivalence holds
in \(\widehat{GT}_0\) or \(\widehat{GT}\)
when \(\alpha\) belongs to the respective group.

For \(\alpha\in\widehat{GT}\), full conjugacy
would in particular imply
\[
                         d_\alpha^2=1+24c_2(u).
 \tag{A.45.3}
\]
Its dyadic component already holds by A.42.
The other components of (A.45.3), and the full
equation (A.45.2), are not proved here.

\begin{proof}
The quotient is obtained by imposing
\(x^2=y^2=1\). Its rotation subgroup is the
full procyclic group generated by \(t=xy\),
and its commutator subgroup is
\(2\widehat{\mathbb Z}\) inside that subgroup.
Since \(u,f\in F'\), write
\(q(u)=t^k\), \(q(f)=t^\ell\) with
\(k,\ell\in2\widehat{\mathbb Z}\).

The defining kernel is preserved by
\(\alpha\), which sends \(a\) to itself
and \(t\) to \(t^{1+2\ell}\).
Its square is one, so \((1+2\ell)^2=1\).
At every odd prime the root is positive,
because A.44 proves identity already modulo
that prime and the only roots of one in
\(\mathbb Z_p\) are \(1,-1\).
At two, \(1+2\ell\equiv1\pmod4\)
also selects the positive root. Thus
\(\ell=0\), and the entire induced action
is the identity. A.44 makes \(2k+1\)
nonzero modulo every odd prime; it is odd
at two. This proves all assertions about
\(d_\alpha\).

For any full unit \(\mu\), the normalized
substitution with parameter \(u\) sends the
two dihedral generators to
\(a,at^{2k+1}\), since a unit power of a
reflection is that same reflection. Their
product is \(t^{d_\alpha}\). The required
third-cusp power is \(t^\mu\). Rotations in
\(D_\infty^{\wedge}\) are conjugate exactly
when their exponents agree or are opposite:
its rotation subgroup is abelian and the
other coset acts by one global inversion.
Hence any full solution has
\(\mu=d_\alpha\) or \(-d_\alpha\).

The third-cusp equations for opposite signs
are equivalent, because the product of the
two inverses is conjugate to the inverse
of the product. A.39 now proves (A.45.2)
as the exact full conjugacy criterion,
including its two membership descents.
The necessary coefficient relation A.4
gives (A.45.3) when a solution exists.

Finally, for a full \(\widehat{GT}\)-involution,
A.42 supplies a prescribed peripheral
conjugator in the maximal pro-two quotient.
Project it to the pro-two dihedral quotient
\(\mathbb Z_2\rtimes C_2\). The same rotation
calculation gives \(\mu_2=\pm(d_\alpha)_2\),
while A.38 gives \(\mu_2^2=1+24c_2(u)_2\).
This proves the stated dyadic component.
\end{proof}

A.45 removes the variable full multiplier
from A.39's finite-quotient test: one may
test the fixed unit \(d_\alpha\) in every
characteristic finite quotient. It does not
replace that full conjugacy test by maximal
pro-prime tests. The comparison between its
dihedral value and A.42's other local roots
is a remaining concrete compatibility issue.

\statement{A.46}{Proved}{Generation over the full elementary abelian two quotient}
Let \(\alpha,u,\psi\) be as in A.44, and put
\[
 V_4=\ker\bigl(F\longrightarrow C_2^2\bigr)
        =\overline{F'F^2}.
\]
For every prime \(p\), the canonical substitution
\(\psi(x)=x\), \(\psi(y)=u^{-1}yu\) induces a
continuous automorphism of
\[
                         F/O^p(V_4).
 \tag{A.46.1}
\]
It intertwines the induced involution with the standard
negative involution. This asserts generation and unmarked
relative conjugacy; it does not assert that this automorphism
preserves the third cusp.

\begin{proof}
The substitution is the identity on \(F^{\mathrm{ab}}\),
so preserves \(V_4\), and then preserves its fully invariant
pro-\(p\) residual. Set \(Q=F/V_4=C_2^2\).
For odd \(p\), the Frattini module
\[
 M=V_4/\overline{V_4'V_4^p}
\]
is the homology of the finite \(Q\)-cover of the
bouquet of two circles. Its cellular complex in degrees
one and zero is
\[
 \mathbb F_p[Q]^2\longrightarrow\mathbb F_p[Q],
 \qquad (a,b)\longmapsto
       a(\bar x-1)+b(\bar y-1).
\]
All four characters of \(Q\) take values in
\(\mathbb F_p\). On the trivial character the displayed
differential is zero, and on each nontrivial character
it has rank one. Consequently
\[
             \dim M^{\mathrm{triv}}=2,
 \qquad \dim M^\chi=1\quad(\chi\ne1).
 \tag{A.46.2}
\]
The induced map \(\psi_M\) is \(Q\)-equivariant,
because \(\psi\) induces the identity on \(Q\).
The inclusion \(V_4\hookrightarrow F\) identifies
the trivial summand with \(H_1(F,\mathbb F_p)\):
transfer followed by inclusion multiplies by four,
and the two spaces both have dimension two.
Naturality and \(\psi_{\mathrm{ab}}=1\) therefore
make \(\psi_M\) the identity on this summand.

Each remaining one-dimensional summand is detected by
an explicit finite dihedral quotient. Write
\[
 D_{2p}=\langle a,t\mid a^2=t^p=1,
                          \ ata=t^{-1}\rangle.
\]
For the character \(\chi(\bar x)=\chi(\bar y)=-1\),
use \(x\mapsto a,y\mapsto at\). The image of
\(V_4\) is the nontrivial rotation subgroup,
with precisely this character for conjugation.
A.44 says that its composite with \(\psi\)
is still onto, so \(\psi_M\) is nonzero on this
character line.

For \((\chi(\bar x),\chi(\bar y))=(-1,+1)\),
use \(x\mapsto a,y\mapsto t\). Since
\(u\in F'\), its image is a rotation; it commutes
with \(t\). The composite with \(\psi\) therefore
retains both generators \(a,t\). For the remaining
character \((+1,-1)\), use
\(x\mapsto t,y\mapsto a\); the composite retains
\(t\) and a conjugate reflection, and again is
onto. In each case \(V_4\) maps onto the rotation
subgroup, since the square of its rotation generator
is nontrivial for odd \(p\). Equivariance shows that
this nonzero map on \(M\) factors through the
indicated character line. Thus \(\psi_M\) is
nonzero on every one-dimensional summand, and is
an automorphism of \(M\).

The induced endomorphism of the free pro-\(p\)
group \(V_4^{(p)}\) is now onto by the Frattini
criterion. Together with the identity on \(Q\),
this proves surjectivity of (A.46.1). A finitely
generated profinite group is Hopfian, so the induced
endomorphism is an automorphism.

At \(p=2\), the group (A.46.1) is exactly the
maximal pro-two quotient of \(F\): it is an
extension of a pro-two group by \(C_2^2\), and
every pro-two quotient of \(F\) restricts to
a pro-two quotient of \(V_4\). Its Frattini map
is the identity, which proves the assertion there.
Finally \(\alpha\psi=\psi c\) is already an
identity on the full group, so descends to every
quotient just considered.
\end{proof}

\statement{A.47}{Proved}{A genuine noncentral two layer in the compatible parity-one tower}
Let \(V\) be the augmentation submodule of the natural
permutation module \(\mathbb F_2^9\) for \(A_9\), and put
\(S=V\rtimes A_9\). Enlarge the tower of A.41 by all
tripod quotients of \(\widehat K\) onto \(S\) whose
three retained cusps form an antipodal pair and one
other vertex and have images of order seven. This
family is nonempty. Its actual joint quotient with
\(C=\Pi/R\), and its actual scalar central refinement,
carry compatible negative involutory images with
\(\Delta=1\). One full negative element realizes these
images simultaneously with the whole tower of A.41.

The same assertion holds for all such regular-head
tripod families with targets \(W\rtimes T\), where
\(T\) is a finite nonabelian simple group excluded
from A.31's congruence composition factors and \(W\)
is a nontrivial finite irreducible \(T\)-module.
No nonemptiness assertion is made for every possible
module. The explicit group \(2^8\rtimes A_9\) proves
that the enlarged tower contains a noncentral
chief layer of characteristic two. No full
involution or cofinality assertion follows.

\begin{proof}
\emph{The concrete noncentral quotient.}\par
The module \(V\) has dimension eight. For every
nonzero even-weight vector, a suitable three-cycle
changes it in exactly two positions. The orbit of
a resulting pair-vector spans \(V\), so \(V\)
is simple. It is nontrivial, and therefore has
neither invariants nor a trivial quotient. Thus
\(S\) is perfect and centerless.

Since nine is odd, \(\mathbb F_2^9\) splits as the
trivial module plus \(V\). Eckmann--Shapiro gives
\[
 H^1(A_9,\mathbb F_2^9)
       =H^1(A_8,\mathbb F_2)=0,
 \qquad H^1(A_9,V)=0.
\]
Here both alternating groups are perfect; the
cohomological comparison is
\citet[Proposition~3.3]{ABrown2010}.
Use the explicit triple from B.57,
\[
 X=(1234567),\qquad Y=(1254389),
 \qquad Z=(XY)^{-1},
\]
and its affine lifts
\[
 \widetilde X=(0,X),\qquad
 \widetilde Y=(e_1+e_2,Y),\qquad
 \widetilde Z=(e_1+e_9,Z).
 \tag{A.47.1}
\]
They multiply to one and all have order seven.
The exact certificate
\path{check_A47_affine_A9.py}
checks these identities, all 256 simultaneous
vector conjugators, and the generation of the
181440-element simple head. The two translation
coordinates in (A.47.1) are not a simultaneous
coboundary \(((1-X)v,(1-Y)v)\).

The group generated by the first two lifts projects
onto \(A_9\), and its intersection with \(V\) is
a submodule. If that intersection were zero,
it would be a complement, hence conjugate to the
standard complement by \(H^1(A_9,V)=0\). This is
exactly the excluded coboundary condition. Its
intersection is therefore \(V\), proving generation
of \(S\). The script additionally counts 1024
order-seven lift pairs and 256 simultaneous
coboundaries. It does not enumerate \(S\) itself;
the preceding argument proves its generation.

\emph{Invariant joint bundles.}\par
Take all kernels in the stated finite family.
They are permuted by \(\widehat{GT}_0\): cusp
supports and the antipodal relation are preserved,
and unit powers preserve order seven. Group the
kernels by their induced simple head kernel
\(\widehat K\to A_9\). For a fixed head, their
joint image \(P\) is a subgroup of
\(V^n\rtimes A_9\) surjecting onto \(A_9\).
Identifications of the head may twist the module;
for \(A_9\) those twists preserve the natural
module, and more generally they remain nontrivial
simple modules. Its kernel \(N\) is a submodule
of a finite semisimple sum of nontrivial modules.
Consequently
\[
                 [N,A_9]=N,
 \qquad N^{A_9}=0.
\]
It follows that \(P\) is perfect and centerless.
No splitting of \(P\) is needed.

Distinct head kernels give independent bundles.
A proper subdirect product would, by finite
Goursat, give a common nontrivial quotient. That
quotient is perfect and has a nonabelian simple
quotient. A homomorphism from a product to a
nonabelian simple group factors through one
factor, since the factor images commute and are
normal. Thus such a common quotient would identify
two head kernels. Induction proves the full product.

Every head has three nonzero cusp supports. Its
\(A_5\)-stabilizer is trivial by the antipodal-triple
calculation in B.57. Inducing one whole bundle
per head orbit therefore gives the actual finite
quotient
\[
              q:\Pi\twoheadrightarrow M\rtimes A_5,
\]
where \(M\) is a product of finite perfect
centerless bundles, regularly permuted on each
\(A_5\)-orbit. Its kernel is precisely the
intersection of the invariant family.
The quotient \(M/q(R)\) is a finite quotient of
the open normal subgroup \(\widehat K/R\) of
\(C\), so contains no \(A_9\) composition factor.
It is perfect, and its remaining composition
factors are cyclic of order two. It is therefore
trivial. Hence \(q(R)=M\), and the actual joint
quotient is \(M\rtimes C\).

\emph{The actual lift.}\par
The proof of A.40 uses only finite perfectness
of \(M\) and the regular permutation of its
bundles. Its universal central cover is finite
by the argument in A.41. For
\(W=R\cap\ker q\), the actual refinement
\(\Pi/[W,\Pi]\) has finite perfect kernel
\(M^\#\); the regular two- and three-cycles
allow its actual torsion generators to be
conjugated into the pure \(C\)-section.
Thus that proof gives
\[
 \Pi/[W,\Pi]
       =M^\#\rtimes C
       =(M^\#\rtimes A_5)\times_{A_5}C.
\]
The automorphism \(\operatorname{id}_{M^\#}\rtimes h_C\)
fixes the actual cusp and retains the full
\(C\)-frames. A.3 gives the full ordinary lift,
whose product with \(c\) has the required
negative involutory image. The proof of A.43
also applies: the bundle is centerless and
perfect, and the remaining action on its two
central subquotients factors through the perfect
group \(A_5\). Thus
\( [[W,\Pi],\Pi]=[W,\Pi]\), and the actual
first central refinement already has the
compatible order-two action.

For a general finite union of targets
\(W\rtimes T\), the kernels over a fixed simple
head lie in a semisimple sum of nontrivial
irreducible modules, possibly in different
characteristics. The same perfectness,
centerlessness, Goursat, and regular-orbit
arguments apply. A perfect quotient with only
cyclic composition factors is trivial, so the
composition-factor argument still gives
\(q(R)=M\). Joining A.41's simple families only
adds new heads or redundant simple quotients
of existing bundles. Finite unions therefore
satisfy all the same conclusions. The closed
solution sets retain exactly the prescribed
\(C\)-action and frames and the identity on
each finite quotient. Compactness, as in A.41,
provides one full element realizing their
entire compatible tower. Its full square is
still not determined.
\end{proof}

\statement{A.48}{Proved}{Nonregular orbit-one simple quotients also admit compatible images}
The parity-one tower of A.47 can be enlarged by
all finite invariant families of simple quotients
of \(\Pi=C_2\amalg C_3\) excluded from the
congruence composition factors of A.31, together
with their actual scalar central refinements.
These quotients have \(A_5\)-orbit length one
on their simple \(\widehat K\)-head kernels;
regular permutation of head factors is not
required for this assertion. One explicit
nonempty family consists of all epimorphism
kernels \(\Pi\to A_9\) for which \(\tau\)
has order fifteen.

One full negative element still realizes the
compatible involutory images with \(\Delta=1\).
This statement does not provide a full
involution or control a nonperfect noncentral
kernel over a mixed solvable head.

\begin{proof}
\emph{A central-extension lemma.}\par
Let \(R=[R,\Pi]\), \(C=\Pi/R\), and let
\(h_0\) be a full positive automorphism
preserving \(R\) and fixing \(\tau\) literally.
Write \(h_C\) for its induced action and retain
its full \(s/r\) conjugators. Consider a central
extension
\[
                 1\longrightarrow A\longrightarrow
                 C^\#\longrightarrow C\longrightarrow1
\]
with \(A\) finite abelian. Choose arbitrary
lifts \(b_s,b_r\) of the named elements of
\(C\). We claim that \(A\) injects into
\((C^\#)^{\mathrm{ab}}\), and that \(h_C\)
has a lift fixing \(A\) and \(b_sb_r\), with
the prescribed \(C\)-frames lifted.

Pull back to \(P=C^\#\times_C\Pi\). Its
normal subgroup
\[
                  R_0=\{(1,r):r\in R\}
\]
affords the same conjugation action as \(R\)
in \(\Pi\), so \(R_0=[R_0,P]\). Hence
\(P^{\mathrm{ab}}=(C^\#)^{\mathrm{ab}}\).
Set \(S=(b_s,s)\), \(T=(b_r,r)\),
\(a=b_s^2\), and \(b=b_r^3\). The central
extension of the free profinite product has
the presentation
\[
 P=\langle A\text{ central},S,T\mid
                  S^2=a,\ T^3=b\rangle.
 \tag{A.48.1}
\]
Indeed the amalgamated product of the two finite
central extensions over \(A\) has quotient
\(\Pi\); its natural map to \(P\) is the
identity on both \(A\) and that quotient,
so is an isomorphism.

The relations in the abelianization do not
kill a nonzero element of \(A\). To see this
without an integral-presentation assumption,
embed \(A\) in a divisible torsion abelian
group and choose a half of \(a\) and a third
of \(b\). They and \(A\) generate a finite
abelian group to which (A.48.1) maps, injectively
on \(A\). This proves the claimed injection.

Lift the two full \(h_0\)-conjugators from
\(\Pi\) to \(P\), conjugate \(S,T\) by
those lifts, and fix \(A\). Presentation
(A.48.1) gives an automorphism \(\phi\) of
\(P\): it covers the automorphism \(h_0\)
and is the identity on its central kernel.
For \(r\in R_0\), its image differs from
the canonical element of \(R_0\) over
\(h_0(r)\) by a central character
\(\chi:R\to A\). Equivariance makes
\(\chi\) invariant under \(\Pi\), so
\(R=[R,\Pi]\) kills it. Thus \(\phi\)
descends to an automorphism \(k\) of
\(C^\#\), fixing \(A\) and retaining the
specified \(C\)-frames.

Finally \(\phi(ST)(ST)^{-1}\in A\), since
\(h_0(\tau)=\tau\). The automorphism \(\phi\)
is the identity on \(P^{\mathrm{ab}}\):
it fixes \(A\) and conjugates \(S,T\).
Injectivity of \(A\) into this abelianization
forces that difference to be one. Hence
\(k(b_sb_r)=b_sb_r\) literally. No splitting
of \(C^\#\to C\) has been assumed.

\emph{The nonregular quotient.}\par
There is an actual epimorphism \(\Pi\to A_9\)
with
\[
 \begin{aligned}
 s&=(12)(34)(56)(78),\\
 r&=(138)(249)(576),\\
 \tau&=(149)(23758).
 \end{aligned}
 \tag{A.48.2}
\]
The exact certificate
\path{check_A48_orbit_one.py}
checks the orders and enumerates the generated
181440-element group. Its \(\widehat K\)
image is a nontrivial normal subgroup of
\(A_9\), hence all of \(A_9\); otherwise
\(A_9\) would be a quotient of \(A_5\).
The kernel on \(\widehat K\) is \(\Pi\)-normal,
so has orbit length one. Each of the twelve
cusps maps to a conjugate of \(\tau^5\),
of order three, distinguishing this head from
the tripod heads of A.47.

Take all \(\Pi\to A_9\) kernels with cusp
order fifteen. The finite family is
\(\widehat{GT}_0\)-invariant, and its joint
simple image is \(T=A_9^n\). The kernel
\(R\) maps onto \(T\), because its quotient
\(C\) has no \(A_9\) composition factor.
For \(W=R\cap\ker(\Pi\to T)\), the actual
joint quotient is therefore
\[
                             \Pi/W=T\times C.
\]
The same construction works for any finite
invariant family of the simple targets in
the statement.

\emph{The actual central product.}\par
Put \(E=\Pi/[W,\Pi]\) and
\(A=W/[W,\Pi]\). As in A.40, relative
perfectness makes
\(M=R/[W,\Pi]\) the finite perfect image
of the universal central cover of \(T\).
Let \(C^\#\) be the preimage of the pure
\(C\)-section in \(T\times C\).
An element of \(C^\#\) acts trivially on
\(T\), hence on its perfect central cover
\(M\): the difference would be a homomorphism
from \(M\) to the abelian central kernel.
Consequently
\[
                 E=MC^\#,
 \qquad [M,C^\#]=1,
 \qquad M\cap C^\#=A.
 \tag{A.48.3}
\]
Choose \(m_s,m_r\in M\) lifting the actual
finite components and write
\(s=m_sb_s\), \(r=m_rb_r\) with
\(b_s,b_r\in C^\#\). Then
\(\tau=m_sm_rb_sb_r\).
The lemma, applied to the full \(h_0\) already
provided by A.33, gives \(k\) fixing
\(A\) and \(b_sb_r\). Thus
\(\psi_E=\operatorname{id}_M\,k\) is a
well-defined automorphism of (A.48.3).
It fixes the actual cusp. Its prescribed
conjugators lie in \(C^\#\) and commute
with \(m_s,m_r\), so retain exactly the
required \(C\)-frames on the actual
\(s,r\). A.3 therefore gives the full
ordinary positive lift.

Genuine \(c\) preserves \(E,M,C^\#\).
The negative map \(b=\psi_Ec_E\) has square
one on \(M\) and on \(C\), and fixes \(A\)
after squaring. Its square on \(C^\#\)
is therefore a central character \(C\to A\).
It fixes \(\tau\) and \(M\), hence fixes
\(b_sb_r\). Since \(C^{\mathrm{ab}}=C_6\)
is generated by \(\tau_C\), the character
is zero. Thus \(b^2=1\) already on the
actual refinement \(E\).

There is no additional scalar central layer.
Writing \(W_1=[W,\Pi]\) and
\(W_2=[W_1,\Pi]\), the group \(R/W_2\)
is a finite perfect central cover of \(M\).
The preimage of \(A\) in it is central,
by perfectness. The preimage of \(C^\#\)
also centralizes that perfect cover, since
it acts trivially on \(M\). These subgroups
generate \(\Pi/W_2\), so \(W/W_2\) is
central and \(W_1=W_2\).

\emph{Compatibility with the earlier tower.}\par
For a finite regular family of A.47, its
\(\widehat K\)-head supports have size three;
the orbit-one heads have all twelve supports.
They are therefore independent. In particular
\(q(W_{\mathrm{reg}})=T\), and
\[
 q([W_{\mathrm{reg}},\Pi])=[T,T]=T.
\]
The stabilized actual refinement of the
regular family may consequently be used
as the new base in the lemma: its residual
kernel is still relative-perfect and
surjects onto \(T\). No composition-factor
exclusion for this enlarged base is needed.
The argument then retains both finite
requirements and their literal frames.
Finite unions and the same closed-set
compactness argument as A.41 give the
claimed single full element. Only its
specified quotient images have been
proved involutory.
\end{proof}

\statement{A.49}{Proved}{A nonperfect mixed-head layer has an exact Clifford refinement}
Let \(H=\operatorname{AGL}_1(\mathbb F_7)\), let
\(V\subset\mathbb F_2^7\) be its augmentation module,
and put \(G=V\rtimes H\). The complete family of
\(\Pi\to G\) epimorphism kernels consists of two
kernels, exchanged by \(c\). Write their joint
quotient as
\[
 \overline G=\overline V\rtimes\overline H,
 \qquad \overline V=V_+\oplus V_-,
 \qquad \overline H=C_7^2\rtimes C_6.
\]
For \(R,C\) as in A.31--A.34 and
\(W=R\cap\ker(\Pi\to\overline G)\), one has
\[
 q(R)=\overline V,
 \qquad
 \Pi/[W,\Pi]
   \simeq\bigl((P_+\times P_-)\rtimes\overline H\bigr)
                       \times_{\overline H}C,
 \tag{A.49.1}
\]
where \(P_\pm\simeq2_+^{1+6}\) are the specified
Clifford central covers of \(V_\pm\). Moreover
\[
                       [[W,\Pi],\Pi]=[W,\Pi].
 \tag{A.49.2}
\]
A full negative element with \(\Delta=1\) realizes
an involutory image on this actual refinement,
compatibly with the entire earlier tower of A.48.
The residual image \(\overline V\) is noncentral
and nonperfect. This is a particular mixed
solvable-head calculation, not a general norm
vanishing or a full order-two construction.

\begin{proof}
\emph{The finite heads and their cohomology.}\par
Label the seven points by \(\mathbb F_7\). There
is an actual epimorphism \(\Pi\to G\) with
\[
 s(t)=-t,\quad r(t)=2t+1,
 \qquad \widetilde s=(0,s),
 \quad \widetilde r=(e_2+e_5,r).
 \tag{A.49.3}
\]
The certificate \path{check_mixed_affine7.py}
verifies the lift orders two and three, enumerates
all \(2688=2^6\cdot42\) generated elements,
and verifies cusp order twelve.

The module \(V\) is irreducible. Under \(C_7\)
it has two irreducible three-dimensional
constituents, interchanged by a primitive
\(C_6\)-element. Its dot product is alternating
and nondegenerate, so it is self-dual. Also
\[
                         H^1(H,V)=0.
 \tag{A.49.4}
\]
Indeed \(\mathbb F_2^7=\mathbb F_2\oplus V\),
and Shapiro identifies its first cohomology
with \(H^1(C_6,\mathbb F_2)\). The trivial
summand maps isomorphically to this group
by restriction, since \(H^{\mathrm{ab}}=C_6\).
The comparison is the same
\citet[Proposition~3.3]{ABrown2010} used in A.47.

A cocycle of \(\Pi\) with this module is a pair
\((v_s,v_r)\) satisfying the order-two and
order-three norm equations. On the seven points, \(s\) has three transpositions
and one fixed point, while \(r\) has two three-cycles
and one fixed point. The norm kernels on \(V\)
therefore have dimensions three and four. Thus
the solution space has dimension seven; the
coboundaries have dimension six because \(V^H=0\). Both counts are checked exactly
in the certificate. Thus
\[
                         \dim H^1(\Pi,V)=1.
 \tag{A.49.5}
\]
Its nonzero class is (A.49.3). A proper subgroup
of \(G\) surjecting onto \(H\) would intersect
the simple \(V\) trivially, hence be a complement;
(A.49.4) makes it conjugate to the standard one.
This is precisely the zero cohomology class.

The subgroup \(V=O_2(G)\) is characteristic,
since \(H\) has no nontrivial normal two-subgroup.
There are two head kernels \(\Pi\to H\),
according as the multiplier of \(r\) is two
or four. Affine conjugation normalizes the
fixed point of \(s\) to zero and the nonzero
\(r\)-displacement to one. Every automorphism
of \(H\) is inner: its characteristic \(C_7\)
and the defining conjugation relation force
identity on \(C_6\), and the remaining scalar
and translation choices are inner. Equation
(A.49.5) therefore gives just one generating
\(G\)-kernel above each head. This is the full
finite family, so it is \(\widehat{GT}_0\)-invariant.
The formula \(c(r)=sr^{-1}s\), \(c(s)=s\)
exchanges its two members.

Their joint head has order \(294=7^2\cdot6\),
with inverse primitive \(C_6\)-characters on
its two \(C_7\)'s. The two modules \(V_+,V_-\)
are nonisomorphic simples for this head. The
joint image has the entire kernel
\(\overline V\): otherwise one simple summand
would have zero projection from that kernel,
making \(G\), of order 2688, a quotient of
the 294-element head. Semisimplicity of this
direct sum proves the asserted joint quotient.

\emph{The actual residual image.}\par
The congruence-kernel image in \(\overline H\)
is exactly \(C_7^2\). The quotient \(C_6\)
is congruence, whereas a finite solvable
congruence quotient has only two and three
as composition primes. For completeness,
\(\operatorname{SL}_2(\mathbb Z/p^n)\) is
perfect for \(p\geq5\): elementary unipotents
are commutators with a diagonal unit whose
square minus one is invertible, and they
generate the group. The remaining local
factors have only two- and three-primary
composition factors. The two different
\(C_6\)-characters then give both independent
seven-directions.

Normality forces the congruence-kernel image
in \(\overline G\) to be
\(\overline V\rtimes C_7^2\). Its nilpotent
residual is \(\overline V\): the quotient is
nilpotent, while every nilpotent image must
kill \([\overline V,C_7^2]=\overline V\),
since its normal two- and seven-parts commute.
The residual commutes with an epimorphism by
A.34's compact comparison. Hence
\(q(R)=\overline V\), as claimed.

The retained positive \(h_C\) acts identically
on \(\overline H\). At a principal congruence
level divisible by twelve the kernel's image
in \(\overline H\) is a seven-group, so this
quotient factors through the relative pro-seven
level, where A.33's positive action is identity.
The nonzero class in (A.49.5) does not inflate
from \(C\), since it is nonzero on \(R\).
Inflation is injective; therefore
\[
 H^1(C,V_+)=H^1(C,V_-)=0,
 \qquad H^1(C,\overline V^\vee)=0,
 \tag{A.49.6}
\]
where \(\vee\) denotes finite Pontryagin duality.

\emph{The actual central-kernel bound.}\par
Put \(E=\Pi/[W,\Pi]\), \(A=W/[W,\Pi]\),
and \(M=R/[W,\Pi]\). Then \(M\) is a
central extension of \(\overline V\) by
\(A\), and \(M=[M,E]\). Work in an arbitrary
finite quotient retaining \(\overline G\),
with images \(M_i,A_i\). If \(A'_i\) is
the image of \(A_i\) in \(M_i^{\mathrm{ab}}\),
duality gives an exact sequence
\[
 0\longrightarrow\overline V^\vee
  \longrightarrow(M_i^{\mathrm{ab}})^\vee
  \longrightarrow(A'_i)^\vee\longrightarrow0.
\]
The middle invariants vanish by relative
perfectness; \(A'_i\) is a trivial \(C\)-module.
The invariant long exact sequence and
(A.49.6) force \(A'_i=0\). Thus
\(A_i=[M_i,M_i]\).

Its commutator pairing is an alternating,
bilinear, \(C\)-invariant pairing on
\(\overline V\). Hence \(A_i\) is a quotient
of \((\bigwedge^2\overline V)_C\), of
dimension two. Each summand contributes one
direction: sharp two-transitivity gives
endomorphism dimension two for the seven-point
permutation module, and its splitting
\(1\oplus V\) gives
\(\operatorname{End}_H(V)=\mathbb F_2\).
Its nonzero invariant form is the dot product.
The mixed tensor summand has no coinvariants,
since the first \(C_7\) has no invariants on
\(V_+\) and is trivial on \(V_-\).
Consequently
\[
                          |A_i|\leq4,
 \qquad |M_i|\leq2^{14}.
 \tag{A.49.7}
\]
No finiteness of the original \(A\) was assumed.

\emph{The Clifford quotient attains the bound.}\par
Take seven symbols \(e_i\) with \(e_i^2=1\)
and \(e_ie_j=-e_je_i\) for \(i\ne j\).
The even-subset elements, with their central
signs, form a group \(P\) of order 128,
quotient \(V\), and square function
\[
                         q(v)=\operatorname{wt}(v)/2\pmod2.
\]
It has 36 zeros, so \(P\simeq2_+^{1+6}\).
Permuting the seven symbols is an honest
\(S_7\)-action, with the reordering sign included.
Thus \(P\rtimes H\) is an actual central
extension of \(G\).

The lifts in (A.49.3) lift further with their
same orders: \(s\) is pure and the indicated
\(r\)-lift has positive central sign. The
certificate \path{check_mixed_affine7_clifford.py}
checks the action law, generation of all
5376 elements, and cusp order twenty-four.
Pair this map with its precomposition by \(c\).
The resulting finite image is
\[
                     \widehat G_0=(P_+\times P_-)
                                           \rtimes\overline H.
\]
Its \(R\)-image is all \(P_+\times P_-\):
it covers \(\overline V\), and commutators
of lifts supported in one summand supply
its central sign independently. Therefore
\(W\) maps to the two central signs and
\([W,\Pi]\) is killed. The actual \(M\)
thus maps onto a group of order \(2^{14}\).
Bound (A.49.7) forces an isomorphism; any
additional kernel would survive in a finite
quotient retaining this one and violate the
bound. The two actual projections now give
(A.49.1), injectively by this identification
and surjectively by correcting a finite
\(P\)-coordinate with an element of \(R\).

\emph{Frames and the constructed involution.}\par
The map \((\operatorname{id}_{\widehat G_0},h_C)\)
is an automorphism of this actual fiber product
and fixes its actual cusp. The finite \(s\)
components are pure, so their centralizer maps
onto the corresponding \(\overline H\)-centralizer.
The actual order-three \(r\)-lift is
\(P_+\times P_-\)-conjugate to the pure lift
by the coprime complement theorem, giving
the same surjectivity for \(r\). The exact
certificate displays a conjugator in the
plus factor. Every retained literal
\(C\)-frame therefore glues to a finite
centralizing element, exactly as in A.40.

Genuine \(c\) exchanges the two finite
Clifford factors. The corresponding negative
automorphism on \(E\) is
\((c_{\widehat G_0},\beta_C)\), of order two.
At this point A.3 gives a full lift of its
positive companion on \(\Pi/W\); the next
paragraph is needed to retain its entire
specified action on \(E\).

\emph{Scalar stabilization and the full lift.}\par
Write \(W_1=[W,\Pi]\), \(W_2=[W_1,\Pi]\),
\(A_2=W_1/W_2\), \(M_2=R/W_2\), and
\(Z_0=W/W_2\). Again work first in finite
quotients retaining the finite Clifford image.
The same dual argument applied to
\[
 0\longrightarrow\operatorname{im}(A_2)
   \longrightarrow M_2^{\mathrm{ab}}
   \longrightarrow M^{\mathrm{ab}}=\overline V
   \longrightarrow0
\]
gives \(M_2^{\mathrm{ab}}=\overline V\).
The commutator of a lift of the old central
\(A_1=A\) with \(M_2\) is a \(C\)-invariant
pairing \(A_1\otimes\overline V\to A_2\).
Both central factors are trivial modules;
for a fixed first lift the central commutator
is a homomorphism in the second argument,
so factors through the displayed abelianization.
Since \((\overline V)_C=0\), this pairing
vanishes. Thus \(Z_0\) centralizes \(M_2\).
Since the old center is \(M'\) and
\(A_2\subseteq M_2'\), one has \(Z_0=M_2'\).
It is an elementary abelian quotient of
\(\bigwedge^2\overline V\), with action
factoring through \(\overline H\).

Let \(J=C_7^2\rtimes C_3\) be its odd
normal subgroup. The two subquotients
\(A_2,A_1\) of \(Z_0\) are trivial
\(J\)-modules. Maschke makes its action
on \(Z_0\) trivial. But
\[
                    \dim(\bigwedge^2\overline V)_J=2.
\]
In each summand the three pairs of opposite
seven-characters are cycled by \(C_3\),
leaving one direction; independent translations
kill the mixed directions. The certificate
checks alternating-square relation rank 14
using only the translation and \(r\),
without the even generator, and verifies
\(|J|=147\). The quotient \(A_1\) already
has dimension two, so \(A_2=0\).
Passing over finite quotients proves (A.49.2).
Now apply A.3 with \(N=W_1\): its actual
central refinement is itself. The specified
framed automorphism of \(E\) therefore lifts
to a full positive \(h\), and \(hc\) has
the asserted involutory image and \(\Delta=1\).

\emph{Retaining the earlier towers.}\par
A finite requirement from A.47--A.48 has
perfect finite heads, whereas \(\overline G\)
is solvable. They have no nontrivial common
quotient. On \(R\), the earlier finite
perfect image and \(\overline V\) likewise
have no common quotient; hence the earlier
kernel still maps onto \(\overline V\).
Taking its relative commutator preserves
that image, since
\([\overline V,\overline G]=\overline V\).
Its stabilized residual is relative-perfect
by A.43 as applied in A.47 and by A.48.
Using that refinement as the new base,
(A.49.5) again proves (A.49.6), since the
nonzero class still detects the residual.
Every argument above therefore retains
that finite requirement and its frames.
Finite unions and compactness give one
full element for the entire asserted tower.
Its full square is not determined.
\end{proof}
